\documentclass[11pt,oneside]{article}

\usepackage[a4paper,margin=27mm,headheight=15pt]{geometry}
\usepackage[T1]{fontenc}
\usepackage[utf8]{inputenc}
\IfFileExists{libertinus.sty}{\usepackage{libertinus}}{\usepackage{lmodern}}
\usepackage{microtype}
\usepackage{amsmath,amssymb,amsthm,mathtools,mathrsfs}
\usepackage{bm}
\usepackage{booktabs,longtable,array,tabularx}
\usepackage{graphicx}
\usepackage{pgfplots}
\usepackage{xcolor}
\usepackage{enumitem}
\usepackage{fancyhdr}
\usepackage{titlesec}
\usepackage{listings}
\usepackage{xurl}
\usepackage{hyperref}
\usepackage[nameinlink,noabbrev]{cleveref}

\definecolor{linkblue}{RGB}{25,75,135}
\definecolor{shadegray}{RGB}{246,247,249}
\definecolor{rulegray}{RGB}{195,201,208}
\pgfplotsset{compat=1.18}
\hypersetup{
  colorlinks=true,
  linkcolor=linkblue,
  citecolor=linkblue,
  urlcolor=linkblue,
  pdftitle={Fabius integration research frontiers},
  pdfsubject={Mixed-status source notebook and outstanding formalization obligations},
  pdfkeywords={Fabius function, Rvachev up-function, inverse, non-elementary, Toeplitz theorem, saddle expansion, Gaussian contraction}
}

\pagestyle{fancy}
\fancyhf{}
\fancyhead[L]{\small Fabius integration research frontiers}
\fancyhead[R]{\small\nouppercase{\leftmark}}
\fancyfoot[C]{\thepage}
\renewcommand{\headrulewidth}{0.4pt}

\titleformat{\section}{\Large\bfseries}{\thesection}{0.7em}{}
\titleformat{\subsection}{\large\bfseries}{\thesubsection}{0.7em}{}
\titleformat{\subsubsection}{\normalsize\bfseries}{\thesubsubsection}{0.7em}{}
\setcounter{secnumdepth}{3}
\setcounter{tocdepth}{2}

\setlist[itemize]{leftmargin=2em,itemsep=0.25em,topsep=0.4em}
\setlist[enumerate]{leftmargin=2.3em,itemsep=0.25em,topsep=0.4em}
\setlength{\emergencystretch}{2em}

\newtheorem{theorem}{Theorem}[section]
\newtheorem{proposition}[theorem]{Proposition}
\newtheorem{lemma}[theorem]{Lemma}
\newtheorem{corollary}[theorem]{Corollary}
\newtheorem{conjecture}[theorem]{Conjecture}
\theoremstyle{definition}
\newtheorem{definition}[theorem]{Definition}
\newtheorem{algorithm}[theorem]{Algorithm}
\newtheorem{example}[theorem]{Example}
\theoremstyle{remark}
\newtheorem{remark}[theorem]{Remark}
\newtheorem{warning}[theorem]{Warning}

\newcommand{\R}{\mathbb R}
\newcommand{\C}{\mathbb C}
\newcommand{\Q}{\mathbb Q}
\newcommand{\N}{\mathbb N}
\newcommand{\Z}{\mathbb Z}
\newcommand{\Up}{\operatorname{up}}
\newcommand{\sinc}{\operatorname{sinc}}
\newcommand{\wt}{\operatorname{wt}_2}
\newcommand{\e}{\mathrm e}
\newcommand{\ii}{\mathrm i}
\newcommand{\dd}{\,\mathrm d}
\newcommand{\Li}{\operatorname{Li}}
\newcommand{\Var}{\operatorname{Var}}
\newcommand{\E}{\mathbb E}
\newcommand{\Prob}{\mathbb P}
\newcommand{\bigO}{\mathcal O}
\newcommand{\smallo}{o}
\newcommand{\supp}{\operatorname{supp}}
\newcommand{\sgn}{\operatorname{sgn}}
\newcommand{\lcm}{\operatorname{lcm}}
\newcommand{\vtwo}{v_2}
\newcommand{\qbinom}[3]{\genfrac{[}{]}{0pt}{}{#1}{#2}_{#3}}
\newcommand{\repo}{\nolinkurl{Analysis/FabiusFunction}}
\newcommand{\lean}[1]{\texttt{#1}}
\newcommand{\anloc}{\mathcal A}
\newcommand{\Gcal}{\mathcal G}
\newcommand{\Pcal}{\mathcal P}
\newcommand{\Wcal}{\mathcal W}
\newcommand{\Id}{\operatorname{id}}
\newcommand{\interior}{\operatorname{int}}

\newenvironment{proofidea}{\par\smallskip\noindent\textit{Proof architecture.}\ }{\hfill$\square$\par\smallskip}
\newenvironment{boxedremark}{%
  \par\smallskip\noindent\begin{center}\begin{minipage}{0.94\textwidth}%
  \color{black}\hrule\vspace{0.6em}\small
}{%
  \vspace{0.6em}\hrule\end{minipage}\end{center}\smallskip
}

\lstdefinestyle{pseudo}{
 basicstyle=\ttfamily\small,
 backgroundcolor=\color{shadegray},
 frame=single,
 rulecolor=\color{rulegray},
 columns=fullflexible,
 keepspaces=true,
 showstringspaces=false,
 xleftmargin=0.7em,
 xrightmargin=0.7em,
 aboveskip=0.8em,
 belowskip=0.8em
}

\title{\bfseries Fabius Integration Research Frontiers\\[0.35em]
\large A mixed-status source notebook for inversion, probability bridges,\\
discrete engines, and saddle algebra}
\author{A research companion to \emph{The Fabius Function and Rvachev's Up-Function}\\
\small based on Vladimir Reshetnikov's \texttt{ProveIt} formal development}
\date{Repository snapshot: 25 August 2026}

\begin{document}
\maketitle

\begin{abstract}
The long article \emph{The Fabius Function and Rvachev's Up-Function} gives a broad
human-readable account of the formal development in \texttt{ProveIt}.  Its breadth is also
its unavoidable weakness: several new branches of the Lean library postdate the article,
and several generic proof engines are named only in the module guide because they have no
counterpart in the exposition.  This companion supplies those missing parts.

The first part develops the total inverse of the bounded Fabius function, transports
infinite flatness into infinite endpoint steepness, and proves the strong non-elementarity
results: no elementary function can agree with the Fabius function, Rvachev's up-function,
or the signed global extension on any set with nonempty interior in the relevant support.
The second part fills in the measure-theoretic bridges between the product probability law,
its survival function, tilted Laplace moments, finite atomic laws, and half-weighted step
approximants.  The third part presents two reusable discrete engines omitted from the main
article: weighted path sums for triangular recurrences and a finite-row Toeplitz convergence
theorem.  The opening audit also records a newly formalized redundancy in the original
compact-support characterization: positivity of its dilation parameter follows from the
remaining hypotheses.  The final and largest part exposes the formal and analytic machinery behind the
all-orders saddle expansion: parameter-dependent Poincar\'e expansions, recursive
exponential and logarithmic coefficients, exact finite polynomial remainders, Gaussian
contraction, uniform polynomial tails, the closed Stirling description of the saddle jets,
and the explicit second periodic correction.

The purpose is complementarity rather than repetition.  Definitions and headline theorems
already proved in detail in the main article are recalled only when they are needed as input.
Every substantial assertion below is either a direct human-readable rendering of a current
Lean declaration, a proof-level explanation of a generic module that the main article
explicitly omits, or a clearly identified elementary corollary of those declarations.
\end{abstract}

\begin{boxedremark}
\textbf{Formalization boundary.}
This is a research-frontier document.  It preserves the source draft used during integration,
including exact Lean-backed results, proof sketches, ordinary-analysis corollaries, stronger
generalizations, and formulas that still lack exact declarations.  Its presence under
\path{non-formalized-research-frontiers} is intentional: no mathematical assertion in this
notebook should be attributed to the formal library merely because a related definition or
weaker theorem exists.  The primary exposition contains the audited, declaration-backed
subset.  Among the obligations retained here are H\"older consequences beyond the
now-formalized inverse-curvature identities,
generic moment-matrix consequences, arbitrary-family Toeplitz wording, generic Gaussian
coefficient bounds, and the passage from the mass top jet to the logarithmic
highest-derivative law.
\end{boxedremark}

\tableofcontents
\clearpage

\section{Scope, conventions, and the audit boundary}
\label{sec:scope}

Let $F:\R\to[0,1]$ denote the bounded Fabius function, totalized in the usual way by
$F(x)=0$ for $x\le0$ and $F(x)=1$ for $x\ge1$.  Let
\begin{equation}
 \Up(x)=F(1-|x|)
 \label{eq:up-fold-supp}
\end{equation}
be Rvachev's compactly supported up-function, and let $\widetilde F$ denote the signed
global extension used in the Lean development.  On $[0,1]$ the basic identities are
\begin{equation}
 F(1-x)=1-F(x),
 \qquad
 F'(x)=2\Up(2x-1),
 \qquad
 F'(x)=2F(2x)\quad(0<x<\tfrac12).
 \label{eq:basic-inputs}
\end{equation}
These are inputs here, not topics to be reproved.

The repository snapshot underlying this supplement is the head of \texttt{main} on
25 August 2026, commit
\begin{center}
 \texttt{5e80805b1ee92697723949d125aa0d6dbf32f538}.
\end{center}
The main article itself records an earlier pinned snapshot and, near the beginning of its
module guide, explicitly says that four clusters have ``no counterpart anywhere in the
article'': Gaussian contraction and Gaussian tails, the formal saddle exponential and
logarithm algebra, the abstract Toeplitz row lemma, and the generic weighted-path solver.
Those four clusters receive complete mathematical chapters below.  Two additional clusters,
the inverse and non-elementarity developments, are absent because they were added after the
article's pinned snapshot.  The current snapshot also adds a proof that positivity of the
dilation parameter in the original compact-support characterization is redundant.  Finally,
several bridges are present only as one- or two-sentence transitions; they are expanded here
into proof-level arguments.

\subsection{Three levels of coverage}

To keep the audit precise, each topic belongs to one of the following levels.
\begin{longtable}{@{}p{0.19\textwidth}p{0.72\textwidth}@{}}
\toprule
Status & Meaning in this companion\\
\midrule
\endhead
\textbf{New branch} & A current Lean module whose mathematical content does not occur in the
main article: principally \path{FabiusInverse.lean}, \path{ElementaryFunction.lean},
\path{NotElementary.lean}, and \path{ProbabilityLaplaceMoments.lean}.\\
\textbf{Missing engine} & A generic theorem family that the main article deliberately leaves in
the module guide: path sums, Toeplitz rows, formal coefficient recurrences, finite
remainders, and Gaussian polynomial estimates.\\
\textbf{Compressed bridge} & A statement whose endpoints occur in the main article but whose
intermediate measure theory, topology, or asymptotic algebra is suppressed.\\
\bottomrule
\end{longtable}

\subsection{Notation that must not collide}

The exact half moments are denoted by $d_n$, following the main article.  The bounded saddle
exponent jets are instead denoted by
\begin{equation}
 Z_n(t),
 \label{eq:Z-notation}
\end{equation}
although their Lean declaration is
\lean{negativeLaplaceBoundedExponentJet}.  This avoids using the same letter for two wholly
different sequences.  We write $L=\log 2$, $\Psi$ for the smooth one-periodic correction in
the negative-Laplace exponent, and $\lambda(x)$ for the exact lower-Lambert phase at small
positive $x$.

\begin{boxedremark}
\textbf{Formal boundary convention.}  When a statement is a theorem in the current Lean
library, the corresponding declaration is named in the margin of the discussion or in the
concordance in \cref{app:concordance}.  When a statement is an immediate classical calculus
corollary but is not presently packaged as a declaration, it is called a \emph{corollary in
ordinary analysis}.  The inverse second-derivative formula and its strict closed-half shape
are now packaged formally; the distinction remains relevant for general higher inverse
derivatives and for the transfer of the highest-jet coefficient from mass coefficients to
logarithmic coefficients.
\end{boxedremark}

\subsection{The positivity hypothesis in the original characterization is redundant}
\label{sec:scale-positive}

The compact-support characterization used in the first Arias de Reyna paper
\cite{arias05442} assumes a smooth function $\phi:\R\to\R$, a real dilation
parameter $k$, and
\begin{equation}
 \operatorname{tsupp}\phi=[-1,1],\qquad
 \phi(x)>0\quad(-1<x<1),\qquad
 \phi(0)=1,
 \label{eq:original-characterization-data}
\end{equation}
together with the derivative law
\begin{equation}
 \phi'(x)=k\bigl(\phi(2x+1)-\phi(2x-1)\bigr).
 \label{eq:original-characterization-derivative}
\end{equation}
The paper states $k>0$ as a hypothesis.  In the current formal development that field is retained
for source fidelity, but a smart constructor proves that its positivity follows from the other
assumptions.

\begin{theorem}[Redundancy of scale positivity]
\label{thm:scale-positive-redundant}
Suppose $\phi\in C^\infty(\R)$ satisfies
\eqref{eq:original-characterization-data} and
\eqref{eq:original-characterization-derivative} for some $k\in\R$.  Then $k>0$.
Consequently these data determine an instance of the full original-characterization predicate
without a separate positivity proof.
\end{theorem}

\begin{proof}
The support identity and continuity imply $\phi(-1)=0$: the function vanishes on
$(-\infty,-1)$, and $-1$ is a limit point of that interval.  The mean-value theorem on
$[-1,0]$ therefore gives a point $c\in(-1,0)$ such that
\begin{equation}
 \phi'(c)=\frac{\phi(0)-\phi(-1)}{0-(-1)}=1.
 \label{eq:original-characterization-mvt}
\end{equation}
Now $2c-1<-1$, so the support condition gives $\phi(2c-1)=0$, whereas
$-1<2c+1<1$, so interior positivity gives $\phi(2c+1)>0$.  Evaluating
\eqref{eq:original-characterization-derivative} at $c$ yields
\begin{equation}
 1=k\phi(2c+1).
 \label{eq:original-characterization-positive-product}
\end{equation}
The second factor is positive, hence $k>0$.
\end{proof}

This is \lean{IsOriginalFabius.mk\_of\_derivative\_law}.  The proof is small, but it clarifies
the logical content of the original theorem: the sign of the scale is forced before the later
Fourier argument identifies its exact value $k=2$.

\part{The inverse and the elementary-function obstruction}

\section{The total inverse of the bounded Fabius function}
\label{sec:inverse}

The main article proves that $F$ is a continuous strictly increasing bijection of $[0,1]$
onto itself.  The inverse development turns that qualitative fact into a reusable global
object whose endpoint behavior can be estimated exactly.

\subsection{Order-isomorphism construction and totalization}

Let $F_{[0,1]}$ be the restriction of $F$ to the subtype $[0,1]$.  Strict monotonicity and
surjectivity package it as an order isomorphism
\begin{equation}
 \mathfrak F:[0,1]\simeq_o[0,1].
 \label{eq:fabius-order-iso}
\end{equation}
Define the clamping map
\begin{equation}
 \pi(y)=\min(1,\max(0,y))
 \label{eq:clamp}
\end{equation}
and the total inverse
\begin{equation}
 I(y)=\operatorname{val}\bigl(\mathfrak F^{-1}(\pi(y))\bigr).
 \label{eq:total-inverse}
\end{equation}
The declarations \lean{fabiusOrderIso} and \lean{fabiusInv} formalize this construction.
Lean uses \lean{Set.IccExtend} rather than spelling out \eqref{eq:clamp}.

\begin{theorem}[Global inverse package]
\label{thm:global-inverse}
The function $I:\R\to\R$ has the following properties.
\begin{enumerate}
 \item $0\le I(y)\le1$ for every $y\in\R$.
 \item $I$ is continuous and nondecreasing on $\R$.
 \item $I$ is strictly increasing on $[0,1]$ and maps that interval bijectively onto itself.
 \item For $x,y\in[0,1]$,
 \begin{equation}
  F(I(y))=y,
  \qquad
  I(F(x))=x.
  \label{eq:two-sided-inverse}
 \end{equation}
 \item The clamped tails are
 \begin{equation}
  I(y)=0\quad(y\le0),
  \qquad
  I(y)=1\quad(y\ge1).
  \label{eq:inverse-tails}
 \end{equation}
\end{enumerate}
\end{theorem}

\begin{proof}
The inverse of an order isomorphism is again an order isomorphism, hence continuous for the
order topologies and strictly increasing.  The interval projection $\pi$ is continuous and
nondecreasing, so the composite \eqref{eq:total-inverse} is continuous and nondecreasing on
all of $\R$.  Its values lie in $[0,1]$ because the codomain of $\mathfrak F^{-1}$ is that
subtype.  On $[0,1]$ the projection is the identity, and the two inverse identities are the
usual \emph{apply--inverse--apply} laws for $\mathfrak F$.  If $y\le0$, then $\pi(y)=0$ and
$\mathfrak F^{-1}(0)=0$ because $F(0)=0$; the right tail is identical.
\end{proof}

The comparison laws are often more useful than direct substitution.

\begin{proposition}[Comparison calculus]
\label{prop:inverse-comparison}
For $x,y\in[0,1]$,
\begin{align}
 I(y)\le x &\iff y\le F(x),
 & I(y)<x &\iff y<F(x),                                    \label{eq:I-left-compare}\\
 x\le I(y) &\iff F(x)\le y,
 & x<I(y) &\iff F(x)<y.                                    \label{eq:I-right-compare}
\end{align}
\end{proposition}

\begin{proof}
Apply the strictly increasing map $F$ to the first comparison, or the strictly increasing map
$I$ to the second, and use \eqref{eq:two-sided-inverse}.  The strict versions are the same
argument with $<$ in place of $\le$.
\end{proof}

\subsection{Reflection and exact inverse values}

The reflection symmetry of $F$ survives totalization without a domain restriction.

\begin{theorem}[Global reflection]
\label{thm:inverse-reflection}
For every $y\in\R$,
\begin{equation}
 I(1-y)=1-I(y).
 \label{eq:inverse-reflection}
\end{equation}
In particular
\begin{equation}
 I(0)=0,
 \qquad
 I(1)=1,
 \qquad
 I\!\left(\frac12\right)=\frac12.
 \label{eq:inverse-basic-values}
\end{equation}
\end{theorem}

\begin{proof}
If $y\le0$ or $y\ge1$, the identity follows directly from the two constant tails.  Suppose
$0<y<1$.  Both $I(1-y)$ and $1-I(y)$ lie in $[0,1]$.  Applying $F$ and using the reflection
law for $F$ gives
\[
 F(1-I(y))=1-F(I(y))=1-y=F(I(1-y)).
\]
Injectivity of $F$ on $[0,1]$ proves \eqref{eq:inverse-reflection}.  Substituting
$y=1/2$ then forces the midpoint identity.
\end{proof}

Exact dyadic evaluation of $F$ immediately becomes exact evaluation of $I$ at a dense family
of rational ordinates.  Let $D_{n,a}$ denote the exact bounded dyadic evaluator, so that for
$0\le a\le2^n$,
\begin{equation}
 D_{n,a}=F\!\left(\frac{a}{2^n}\right)\in\Q.
 \label{eq:dyadic-evaluator-input}
\end{equation}
The evaluator is clamped for $a>2^n$.

\begin{theorem}[Inverting the exact dyadic evaluator]
\label{thm:inverse-dyadic-evaluator}
For all $n,a\in\N$,
\begin{equation}
 I(D_{n,a})=\frac{\min(a,2^n)}{2^n}.
 \label{eq:inverse-dyadic-evaluator}
\end{equation}
In particular,
\begin{equation}
 F\!\left(\frac14\right)=\frac5{72}
 \quad\Longrightarrow\quad
 I\!\left(\frac5{72}\right)=\frac14.
 \label{eq:inverse-five-seventy-two}
\end{equation}
\end{theorem}

\begin{proof}
For $a\le2^n$, combine \eqref{eq:dyadic-evaluator-input} with the left inverse identity in
\eqref{eq:two-sided-inverse}.  For $a\ge2^n$, the clamped evaluator equals $1$ and the result
is $I(1)=1$.
\end{proof}

\subsection{Formal interior calculus and curvature}

The inverse is constructed using order topology alone.  Once the positivity theorem for
$F'$ in the open unit interval is imported, the local inverse theorem supplies a sharper
interior picture; the formal development now packages both the reciprocal derivative and
full interior smoothness directly.

\begin{corollary}[Interior derivative of the inverse]
\label{cor:inverse-derivative}
For $0<y<1$,
\begin{equation}
 I'(y)=\frac{1}{F'(I(y))}
      =\frac{1}{2\Up(2I(y)-1)}.
 \label{eq:inverse-derivative}
\end{equation}
Thus $I$ is $C^\infty$ on $(0,1)$.
\end{corollary}

\begin{proof}
The point $x=I(y)$ lies in $(0,1)$, and strict positivity of $F'(x)$ there makes $F$ a local
$C^\infty$ diffeomorphism.  Differentiate $F(I(y))=y$ and substitute the second identity in
\eqref{eq:basic-inputs}.
\end{proof}

Because $F$ is strictly convex on $[0,1/2]$ and strictly concave on $[1/2,1]$, inversion
reverses those two shapes.

\begin{corollary}[Curvature of the inverse; formalized]
\label{cor:inverse-curvature}
The inverse $I$ is strictly concave on $[0,1/2]$ and strictly convex on $[1/2,1]$.  For every
$0<y<1$,
\begin{equation}
 I''(y)=-\frac{F''(I(y))}{F'(I(y))^3}.
 \label{eq:inverse-second-derivative}
\end{equation}
The midpoint is the reflection-symmetric change of curvature.
\end{corollary}

\begin{remark}
The first identity in \cref{cor:inverse-derivative}, its explicit \(\Up\)-form, strict
positivity, and the resulting interior smoothness are the named Lean declarations
\path{fabiusInv_hasDerivAt},
\path{deriv_fabiusInv}, \path{deriv_fabiusInv_eq_inv_two_mul_rvachevUp}, and
\path{deriv_fabiusInv_pos}, together with
\path{fabiusInv_contDiffOn_Ioo}.  The reciprocal-cubic rule in
\cref{cor:inverse-curvature} is \path{deriv_fabiusInv_hasDerivAt} together with
\path{deriv_deriv_fabiusInv}; its two strict closed-half conclusions are
\path{strictConcaveOn_fabiusInv_firstHalf} and
\path{strictConvexOn_fabiusInv_secondHalf}.  These shape theorems use derivatives only on
the interval interiors, so they do not assert finite endpoint derivatives.
\end{remark}

\section{Flatness of \texorpdfstring{$F$}{F} becomes steepness of \texorpdfstring{$I$}{I}}
\label{sec:inverse-steepness}

The inverse is most interesting at the two endpoints, where the classical formula
\eqref{eq:inverse-derivative} ceases to have a finite limit.  The formal development proves
this without differentiating the inverse: exact upper bounds for $F$ are transported through
monotonicity and then solved as inequalities for $I$.

\subsection{Root-free inverse bounds}

For $n\in\N$, the effective flatness theorem for $F$ says that on the dyadic window
$0\le x\le2^{-n}$,
\begin{equation}
 F(x)\le 2^{\binom{n+1}{2}}x^n.
 \label{eq:effective-flatness-input}
\end{equation}
The sharp factorial refinement is
\begin{equation}
 F(x)\le \frac{2^{\binom{n+1}{2}}}{n!}x^n.
 \label{eq:sharp-flatness-input}
\end{equation}
Substitute $x=I(y)$ and use $F(I(y))=y$.

\begin{theorem}[Flatness transported through the inverse]
\label{thm:inverse-root-free}
Let $n\in\N$ and $y\in[0,1]$.  If
\begin{equation}
 2^n I(y)\le1,
 \label{eq:inverse-scale-hypothesis}
\end{equation}
then
\begin{align}
 y&\le 2^{\binom{n+1}{2}}I(y)^n,                            \label{eq:inverse-basic-root-free}\\
 y&\le \frac{2^{\binom{n+1}{2}}}{n!}I(y)^n.                \label{eq:inverse-sharp-root-free}
\end{align}
It is enough to assume the argument-side condition
\begin{equation}
 0\le y\le F(2^{-n}).
 \label{eq:inverse-argument-window}
\end{equation}
\end{theorem}

\begin{proof}
Under \eqref{eq:inverse-scale-hypothesis}, the point $x=I(y)$ belongs to the window of
\eqref{eq:effective-flatness-input} and \eqref{eq:sharp-flatness-input}.  Their left sides
become $F(I(y))=y$, giving the two conclusions.  Under
\eqref{eq:inverse-argument-window}, the comparison calculus
\eqref{eq:I-left-compare} gives $I(y)\le2^{-n}$, which is exactly
\eqref{eq:inverse-scale-hypothesis}.
\end{proof}

If roots are allowed, \eqref{eq:inverse-sharp-root-free} reads
\begin{equation}
 I(y)\ge
 \left(\frac{n!}{2^{\binom{n+1}{2}}}y\right)^{1/n}
 \qquad\bigl(0\le y\le F(2^{-n})\bigr).
 \label{eq:inverse-root-form}
\end{equation}
The Lean statements retain the root-free form because all constants remain exact and no
case split for real powers is required.

\subsection{An effective infinite-slope estimate}

The case $n=2$ already forces the difference quotient to diverge.  Since
$2^{\binom32}=8$ and $F(1/4)=5/72$, \eqref{eq:inverse-basic-root-free} gives
\begin{equation}
 y\le8I(y)^2
 \qquad\left(0\le y\le\frac5{72}\right).
 \label{eq:inverse-square-bound}
\end{equation}

\begin{theorem}[Effective domination of every slope]
\label{thm:inverse-effective-slope}
Let $M\in\R$ and $y>0$.  If
\begin{equation}
 y\le F(1/4)=\frac5{72},
 \qquad
 8M^2y<1,
 \label{eq:effective-slope-hypotheses}
\end{equation}
then
\begin{equation}
 My<I(y).
 \label{eq:effective-slope-conclusion}
\end{equation}
\end{theorem}

\begin{proof}
Assume contrariwise that $I(y)\le My$.  The left side is nonnegative, so squaring preserves
the inequality and gives $I(y)^2\le M^2y^2$.  Combining this with
\eqref{eq:inverse-square-bound} yields
\[
 y\le8M^2y^2.
\]
Since $y>0$, division by $y$ gives $1\le8M^2y$, contradicting
\eqref{eq:effective-slope-hypotheses}.
\end{proof}

No positivity assumption on $M$ is necessary: for $M\le0$ the conclusion is automatic from
$I(y)>0$, and the same algebra proves it uniformly.

\begin{theorem}[Infinite one-sided endpoint slopes]
\label{thm:inverse-infinite-slope}
The endpoint difference quotients diverge:
\begin{align}
 \frac{I(y)}{y}&\longrightarrow+\infty
   &&(y\downarrow0),                                       \label{eq:inverse-slope-zero}\\
 \frac{1-I(y)}{1-y}&\longrightarrow+\infty
   &&(y\uparrow1).                                         \label{eq:inverse-slope-one}
\end{align}
\end{theorem}

\begin{proof}
Fix a target slope $M$.  Replace it by $|M|+1$ and choose $y>0$ below both
$F(1/4)$ and $[8(|M|+1)^2]^{-1}$.  Then \cref{thm:inverse-effective-slope} gives
$(|M|+1)y<I(y)$, hence $M<I(y)/y$.  This is exactly convergence to $+\infty$ in the
order topology.  For the right endpoint, put $z=1-y$ and use
$I(1-z)=1-I(z)$.
\end{proof}

\begin{corollary}[Failure of endpoint Lipschitz and H\"older bounds]
\label{cor:inverse-not-holder}
The inverse is not locally Lipschitz at $0$ or $1$.  More strongly, for every
$\alpha>0$ and every $C>0$, no estimate
\begin{equation}
 I(y)\le Cy^\alpha
 \label{eq:holder-false}
\end{equation}
can hold throughout a punctured right neighborhood of $0$; the reflected statement holds
at $1$.
\end{corollary}

\begin{proof}
The Lipschitz claim follows immediately from \eqref{eq:inverse-slope-zero}.  For the H\"older
claim, choose $n$ with $1/n<\alpha$.  On the shrinking window
$0<y\le F(2^{-n})$, \eqref{eq:inverse-root-form} gives $I(y)\ge c_ny^{1/n}$ with
$c_n>0$.  If \eqref{eq:holder-false} held, then
$c_n\le Cy^{\alpha-1/n}$, whose right side tends to zero as $y\downarrow0$, a
contradiction.
\end{proof}

\begin{boxedremark}
The windows $y\le F(2^{-n})$ shrink super-exponentially as $n$ grows.  One must therefore
not read \eqref{eq:inverse-root-form} as a lower bound uniform in $n$.  Its strength is
quantifier order: for each prescribed root exponent one obtains a sufficiently small window.
That is exactly what is needed to defeat every fixed positive H\"older exponent.
\end{boxedremark}

\section{Elementary functions through their analytic locus}
\label{sec:elementary}

The non-elementarity proof does not attempt differential algebra or Liouville theory.  It
uses a topological-analytic invariant that is both stronger and easier to formalize: every
one-variable elementary function is analytic on a dense open set, whereas the Fabius family
is nonanalytic at every point of the relevant interval.

\subsection{The exact formal class}

\begin{definition}[The class $\mathsf{Elem}$]
\label{def:elementary-class}
Let $\mathsf{Elem}$ be the smallest class of total functions $\R\to\R$ containing all
constant functions and the identity and closed under
\begin{equation}
 +,\quad\cdot,\quad-,\quad(\cdot)^{-1},\quad
 f\mapsto f^r\ (r\in\R),\quad
 \exp,\ \log,\ \sin,\ \cos,\ \arcsin,\ \arctan,
 \label{eq:elementary-generators}
\end{equation}
where every unary operation is applied pointwise.
\end{definition}

This is the inductive predicate \lean{IsElementary}.  Composition is deliberately not a
constructor.

\begin{proposition}[Closure derived from the generators]
\label{prop:elementary-derived}
The class $\mathsf{Elem}$ is closed under composition.  It is consequently closed under
subtraction, division, natural powers, polynomial and rational functions, square roots,
absolute values, tangent, hyperbolic functions, $\arccos$, $\operatorname{arsinh}$, and
signed odd roots.
\end{proposition}

\begin{proof}
Induct on the construction of the outer function $g$.  Precomposing a constant or the
identity with an elementary $f$ is elementary.  Every remaining constructor commutes with
precomposition: for example $(g_1+g_2)\circ f=(g_1\circ f)+(g_2\circ f)$, and similarly for
the unary generators.  This proves closure under composition.  The remaining operations
are identities in the generated algebra.  For example
\[
 |f|=\sqrt{f^2},
 \qquad
 \tan f=\frac{\sin f}{\cos f},
 \qquad
 \sinh f=\frac{\e^f-\e^{-f}}2.
\]
For an odd positive integer $n$, the signed root may be written
\[
 \frac{f}{|f|}|f|^{1/n},
\]
with value zero at $f=0$ under the total reciprocal convention.
\end{proof}

\begin{remark}[Totalization makes the theorem stronger]
\label{rem:elementary-totalization}
Mathlib treats reciprocal, logarithm, real power, and inverse trigonometric functions as
total maps on $\R$; outside their classical domains they receive specified extension values.
Thus \cref{def:elementary-class} is at least as large as the class represented by classical
expressions on open domains.  Proving that $F$ does not belong to this larger class is a
stronger statement, not a domain-convention loophole.
\end{remark}

\subsection{Analytic loci}

For a function $f:\R\to\R$, define
\begin{equation}
 \anloc(f)=\{x\in\R:f\text{ is real analytic at }x\}.
 \label{eq:analytic-locus}
\end{equation}
Analyticity is local, hence $\anloc(f)$ is open for every $f$.  The substantive theorem is
density for elementary functions.

The only difficult induction step is composition with a unary generator that has finitely
many exceptional values.  The following lemma isolates it.

\begin{lemma}[Finite-bad-value composition principle]
\label{lem:finite-bad-composition}
Let $f,g:\R\to\R$.  Suppose that $\anloc(f)$ is dense and that there is a finite set
$B\subset\R$ such that $g$ is analytic at every point of $\R\setminus B$.  Then
$\anloc(g\circ f)$ is dense.
\end{lemma}

\begin{proof}
Let $U$ be a nonempty open interval.  Because $\anloc(f)$ is dense and open, choose
$x_0\in U\cap\anloc(f)$.  We prove that every punctured neighborhood of $x_0$ contains a
point where $g\circ f$ is analytic.

Fix $c\in B$.  The analytic function $f-c$ near $x_0$ has two possibilities.
\begin{enumerate}
 \item It vanishes identically on a neighborhood of $x_0$.  Then $f$ is locally constant
 equal to $c$, so $g\circ f$ is locally constant and therefore analytic at $x_0$, regardless
 of whether $g$ is analytic at $c$.
 \item It is not locally zero.  The isolated-zero theorem gives a punctured neighborhood of
 $x_0$ on which $f(x)\ne c$.
\end{enumerate}
If the first alternative occurs for any $c\in B$, we are done.  Otherwise intersect the
finitely many punctured neighborhoods.  The intersection still contains points arbitrarily
close to $x_0$; intersect once more with the open set $U\cap\anloc(f)$.  At any resulting
point $x$, the value $f(x)$ avoids $B$, so $g$ is analytic at $f(x)$ and $f$ is analytic at
$x$.  The analytic chain rule makes $g\circ f$ analytic at $x$.  Thus every nonempty open
$U$ meets $\anloc(g\circ f)$.
\end{proof}

\begin{theorem}[Dense open analytic locus]
\label{thm:elementary-dense-analytic}
If $f\in\mathsf{Elem}$, then $\anloc(f)$ is dense and open.  Equivalently, the nonanalytic
locus
\begin{equation}
 \R\setminus\anloc(f)
 \label{eq:elementary-exceptional-set}
\end{equation}
is closed and nowhere dense.
\end{theorem}

\begin{proof}
Openness is true for every analytic locus.  Density follows by induction on the elementary
construction.

Constants and the identity are analytic everywhere.  For sums and products, the
intersection of the two dense open analytic loci is dense and each point in the intersection
is analytic for the result.  Negation is immediate.  For a unary generator, apply
\cref{lem:finite-bad-composition} with the following exceptional sets:
\begin{center}
\begin{tabular}{@{}lll@{}}
\toprule
Generator & Exceptional values & Analytic elsewhere\\
\midrule
$t\mapsto t^{-1}$ & $\{0\}$ & yes\\
$t\mapsto t^r$ & $\{0\}$ & yes for every fixed $r\in\R$\\
$\log t$ & $\{0\}$ & yes under the total real-log convention\\
$\arcsin t$ & $\{-1,1\}$ & yes\\
$\exp t,\sin t,\cos t,\arctan t$ & $\varnothing$ & everywhere\\
\bottomrule
\end{tabular}
\end{center}
The complement of a dense open set is closed with empty interior, hence nowhere dense.
\end{proof}

\subsection{The reusable obstruction}

\begin{theorem}[No agreement on an interior]
\label{thm:elementary-obstruction}
Let $h:\R\to\R$ be nonanalytic at every point of $\interior(S)$, where
$S\subset\R$ has nonempty interior.  No elementary function $g$ can agree with $h$ on all of
$S$.
\end{theorem}

\begin{proof}
If $g$ were elementary, \cref{thm:elementary-dense-analytic} would give a point
$x\in\interior(S)$ at which $g$ is analytic.  Equality on $S$ implies equality on a whole
neighborhood of $x$ contained in $S$, so analyticity transfers from $g$ to $h$ at $x$, a
contradiction.
\end{proof}

The theorem is stronger than saying that $h$ itself is not elementary: it rules out any
piecewise repair outside $S$ and any elementary surrogate on a substantial subset.

\section{Strong non-elementarity of the Fabius family}
\label{sec:not-elementary}

The main article proves nowhere analyticity of $F$ on $[0,1]$, of $\Up$ on $[-1,1]$, and of
the signed extension on its active interval.  Combining those results with
\cref{thm:elementary-obstruction} gives the new non-elementarity branch almost for free.

\begin{theorem}[Strong local non-elementarity of $F$]
\label{thm:fabius-strong-non-elementary}
Let $U\subset[0,1]$ have nonempty interior.  There is no $g\in\mathsf{Elem}$ such that
\begin{equation}
 g(x)=F(x)\qquad(x\in U).
 \label{eq:elementary-F-agreement}
\end{equation}
In particular, no elementary function agrees with $F$ on $(0,1)$, and $F$ is not an
elementary function on $\R$.
\end{theorem}

\begin{proof}
The function $F$ is nonanalytic at every point of $[0,1]$, hence at every point of
$\interior(U)$.  Apply \cref{thm:elementary-obstruction}.
\end{proof}

\begin{theorem}[Strong local non-elementarity of $\Up$]
\label{thm:up-strong-non-elementary}
If $U\subset[-1,1]$ has nonempty interior, no elementary $g$ agrees with $\Up$ on $U$.
Consequently $\Up$ is not elementary.
\end{theorem}

\begin{theorem}[Strong local non-elementarity of the signed extension]
\label{thm:global-strong-non-elementary}
If $U\subset[0,2)$ has nonempty interior, no elementary $g$ agrees with the signed global
extension $\widetilde F$ on $U$.  Consequently $\widetilde F$ is not elementary.
\end{theorem}

\begin{proof}[Proof of the last two theorems]
The corresponding nowhere-analyticity theorems supply the hypothesis of
\cref{thm:elementary-obstruction}.  The global conclusions follow by taking a nonempty open
subinterval of the active support and observing that an elementary representation on all of
$\R$ would in particular agree there.
\end{proof}

\begin{remark}[Why compact support alone is not the argument]
A compactly supported nonzero function cannot be globally real analytic, but elementary
functions need not be analytic everywhere: $|x|$, $x^{-1}$, and totalized logarithms already
have exceptional points.  What defeats elementarity is not compact support by itself; it is
nonanalyticity on an entire interval, while elementary exceptional sets are nowhere dense.
\end{remark}

\begin{remark}[Why this is sharper than a differential-algebra slogan]
The conclusion does not depend on choosing a particular syntax for ``closed form'' beyond
the explicit inductive class.  More importantly, it is local: changing the function outside
a small open interval cannot make it elementary on that interval.  Thus the theorem rules
out not only a global elementary formula for $F$, but also an elementary formula valid on
any nontrivial open patch of its natural domain.
\end{remark}

\part{Probability laws and the bridges the headline formulas conceal}

\section{The product probability law in measure-theoretic form}
\label{sec:product-law}

The random-series representation in the main article is concise because it speaks in the
usual probabilistic shorthand.  The formal construction must establish continuity,
measurability, product independence, pushforward identities, and support before it may use
the same shorthand.  This section records that chain explicitly.

\subsection{Sample space and uniform convergence}

Let
\begin{equation}
 \Omega=[0,1]^{\N}
 \label{eq:sample-space}
\end{equation}
with its product topology and product Borel probability measure
\begin{equation}
 \mathbf P=\bigotimes_{n\ge0}\operatorname{Unif}[0,1].
 \label{eq:uniform-product}
\end{equation}
Write $\omega_n$ for the $n$th coordinate and define
\begin{equation}
 X(\omega)=\sum_{n=0}^{\infty}\frac{\omega_n}{2^{n+1}}.
 \label{eq:weighted-coordinate-sum}
\end{equation}
For every $\omega$,
\begin{equation}
 0\le \frac{\omega_n}{2^{n+1}}\le2^{-n-1},
 \qquad
 \sum_{n\ge0}2^{-n-1}=1.
 \label{eq:coordinate-majorant}
\end{equation}
The Weierstrass test therefore gives uniform convergence on $\Omega$.  Every partial sum is
continuous, so $X$ is continuous and hence measurable.  Termwise comparison in
\eqref{eq:coordinate-majorant} gives the pointwise support statement
\begin{equation}
 0\le X(\omega)\le1.
 \label{eq:X-support}
\end{equation}

Let
\begin{equation}
 \mu=\mathbf P\circ X^{-1}
 \label{eq:weighted-sum-law}
\end{equation}
be the pushforward law.  It is a probability measure and
\begin{equation}
 \mu([0,1])=1,
 \qquad
 \mu(\R\setminus[0,1])=0.
 \label{eq:law-support}
\end{equation}
The Lean declarations are \lean{weightedCoordinateSum},
\lean{continuous\_weightedCoordinateSum}, \lean{weightedSumDistribution}, and the support
lemmas surrounding them.

\subsection{Deleting the head coordinate}

Define the left shift
\begin{equation}
 (T\omega)_n=\omega_{n+1}.
 \label{eq:tail-map}
\end{equation}
The product measure is invariant under this deletion in the sense that $T_\#\mathbf P=\mathbf
P$, and the head coordinate $\omega_0$ is independent of the entire tail $T\omega$.  Reindexing
the series gives the exact pointwise split
\begin{equation}
 X(\omega)=\frac{\omega_0+X(T\omega)}2.
 \label{eq:X-head-tail-split}
\end{equation}
Consequently the pair $(\omega_0,X(T\omega))$ has law
\begin{equation}
 \operatorname{Unif}[0,1]\otimes\mu.
 \label{eq:head-tail-law}
\end{equation}
Pushing \eqref{eq:X-head-tail-split} through this joint law gives the self-similarity identity
\begin{equation}
 \mu=\left(\operatorname{Unif}[0,1]\otimes\mu\right)
       \circ\left[(u,x)\mapsto\frac{u+x}{2}\right]^{-1}.
 \label{eq:law-self-similarity}
\end{equation}

\begin{proposition}[The smoothing equation for the CDF]
\label{prop:CDF-smoothing-detailed}
Let $H(y)=\mu(({-\infty},y])$.  Then for every $y\in\R$,
\begin{equation}
 H(y)=\int_0^1 H(2y-u)\dd u.
 \label{eq:CDF-smoothing-detailed}
\end{equation}
For $0\le y\le1/2$ this reduces to
\begin{equation}
 H(y)=\int_0^{2y}H(t)\dd t.
 \label{eq:CDF-left-detailed}
\end{equation}
\end{proposition}

\begin{proof}
Condition on the head coordinate in \eqref{eq:X-head-tail-split}.  For a fixed $u$,
\[
 \Prob\!\left(\frac{u+X'}2\le y\right)=\Prob(X'\le2y-u)=H(2y-u),
\]
where $X'$ has law $\mu$ and is independent of $u$.  Integrate over $u\in[0,1]$.  If
$0\le y\le1/2$, support implies $H(2y-u)=0$ when $u>2y$; substitute $t=2y-u$ and use the
reflection identity established below, or reverse the interval directly.
\end{proof}

\subsection{Coordinate reflection}

Let $R:\Omega\to\Omega$ be
\begin{equation}
 (R\omega)_n=1-\omega_n.
 \label{eq:coordinate-reflection}
\end{equation}
Lebesgue measure on $[0,1]$ is invariant under $u\mapsto1-u$, so the infinite product measure
is invariant under $R$.  Moreover
\begin{equation}
 X(R\omega)=\sum_{n\ge0}\frac{1-\omega_n}{2^{n+1}}
 =1-X(\omega).
 \label{eq:X-reflection}
\end{equation}
Thus the law is reflection invariant:
\begin{equation}
 (x\mapsto1-x)_\#\mu=\mu.
 \label{eq:law-reflection}
\end{equation}
Since the CDF is continuous, this becomes
\begin{equation}
 H(1-y)=1-H(y).
 \label{eq:CDF-reflection}
\end{equation}
The fixed-point uniqueness theorem for the bounded Fabius equation now identifies
\begin{equation}
 H=F
 \quad\text{on all of }\R.
 \label{eq:CDF-is-F}
\end{equation}
The phrase ``on all of $\R$'' includes the two constant tails; it is not a reference to the
signed extension $\widetilde F$.

\begin{theorem}[Probability representation, global form]
\label{thm:probability-representation-global}
For every $y\in\R$,
\begin{equation}
 F(y)=\Prob\!\left(\sum_{n\ge0}\frac{U_n}{2^{n+1}}\le y\right),
 \label{eq:F-as-CDF-global}
\end{equation}
where the $U_n$ are independent uniform variables on $[0,1]$.  Equivalently,
\begin{equation}
 \Up(x)=\mu(({-\infty},1-|x|])
 \label{eq:up-as-CDF-global}
\end{equation}
for every $x\in\R$.
\end{theorem}

\section{The survival-function Fubini identity}
\label{sec:survival-fubini}

A particularly useful bridge, absent from the main article's probability chapter, identifies
integration against the Fabius law with an ordinary integral against $\Up$.  It is the compact
support version of the layer-cake or tail-integration formula.

For $t\ge0$, support and \eqref{eq:CDF-is-F} give
\begin{equation}
 \mu((t,\infty))=1-F(t)=\Up(t).
 \label{eq:survival-is-up}
\end{equation}
The equality remains correct for $t>1$, when both sides are zero.

\begin{theorem}[Expectation from the survival function]
\label{thm:survival-Fubini}
Let $g,g':\R\to\R$ be continuous and suppose
\begin{equation}
 g(x)-g(0)=\int_0^x g'(t)\dd t
 \qquad(0\le x\le1).
 \label{eq:g-FTC-hypothesis}
\end{equation}
Then
\begin{equation}
 \int_{[0,1]}g(x)\dd\mu(x)
 =g(0)+\int_0^1 g'(t)\Up(t)\dd t.
 \label{eq:survival-Fubini}
\end{equation}
\end{theorem}

\begin{proof}
Put
\begin{equation}
 A=\{(t,x)\in[0,1]\times\R:t<x\},
 \qquad
 K(t,x)=\mathbf1_A(t,x)g'(t).
 \label{eq:Fubini-kernel}
\end{equation}
Continuity of $g'$ on the compact interval makes $K$ integrable with respect to
$\operatorname{Leb}_{[0,1]}\otimes\mu$.  Integrating first in $t$ and using
\eqref{eq:g-FTC-hypothesis} gives, for $x\in[0,1]$,
\begin{equation}
 \int_0^1K(t,x)\dd t=\int_0^xg'(t)\dd t=g(x)-g(0).
 \label{eq:Fubini-first-order}
\end{equation}
Integrating first in $x$ gives
\begin{equation}
 \int K(t,x)\dd\mu(x)=g'(t)\mu((t,\infty))=g'(t)\Up(t)
 \label{eq:Fubini-second-order}
\end{equation}
by \eqref{eq:survival-is-up}.  Fubini equates the two iterated integrals.  Since $\mu$ is a
probability measure supported on $[0,1]$, the integral of the constant $g(0)$ is $g(0)$,
and rearrangement yields \eqref{eq:survival-Fubini}.
\end{proof}

\begin{remark}[Endpoint conventions disappear under integration]
The survival function uses the open ray $(t,\infty)$.  The CDF uses $(-\infty,t]$.  Their
complements are exact, and no endpoint atom enters because $\mu$ has a continuous CDF.
Likewise, replacing $(0,1)$ by any of the standard half-open versions changes no Lebesgue
integral.  The formal proof nevertheless records these null-set conversions explicitly.
\end{remark}

\section{Raw, endpoint, and tilted Laplace moments}
\label{sec:probability-laplace}

The main article develops half moments and the negative generating function from the
analytic side.  The new probability--Laplace module proves that those quantities are exactly
the moments of the product law.

\subsection{Three moment families}

For a probability measure $\nu$ supported on $[0,1]$, define
\begin{align}
 U_n(\nu)&=\int_{[0,1]}(1-x)^n\dd\nu(x),                  \label{eq:endpoint-moment-def}\\
 L_k(\nu;s)&=\int_{[0,1]}\e^{-sx}x^k\dd\nu(x),           \label{eq:raw-laplace-def}\\
 J_k(s)&=\int_0^1t^k\Up(t)\e^{-st}\dd t.                 \label{eq:survival-tilt-def}
\end{align}
For real $s$, the integrand in \eqref{eq:raw-laplace-def} is nonnegative, so
\begin{equation}
 L_k(\nu;s)\ge0.
 \label{eq:laplace-moment-nonnegative}
\end{equation}

Apply \cref{thm:survival-Fubini} to $g(x)=\e^{-sx}x^k$.  For $k=0$ one obtains
\begin{equation}
 L_0(\mu;s)=1-sJ_0(s).
 \label{eq:laplace-zero-survival}
\end{equation}
For $k\ge0$, applying it to $g(x)=\e^{-sx}x^{k+1}$ gives
\begin{equation}
 L_{k+1}(\mu;s)=(k+1)J_k(s)-sJ_{k+1}(s).
 \label{eq:laplace-successor-survival}
\end{equation}
These equations are exactly the defining clauses of
\lean{fabiusLaplaceMoment}.

\begin{theorem}[Law moments equal analytic Fabius moments]
\label{thm:law-moments-identification}
For every $k\in\N$ and $s\in\R$,
\begin{equation}
 L_k(\mu;s)=M_k(s),
 \label{eq:law-M-identification}
\end{equation}
where
\begin{align}
 M_0(s)&=G(-s),                                            \label{eq:M-zero}\\
 M_{k+1}(s)&=(k+1)J_k(s)-sJ_{k+1}(s)                      \label{eq:M-successor}
\end{align}
are the tilted Fabius moments and $G$ is the generating function of the main article.
In particular $M_k(s)\ge0$ for real $s$.
\end{theorem}

\begin{proof}
The case $k=0$ is the integration-by-parts identity
\eqref{eq:laplace-zero-survival}, which is also the negative generating-function formula.
The successor case is \eqref{eq:laplace-successor-survival}.
\end{proof}

\subsection{Differential recurrence}

Differentiation under the compact integral in \eqref{eq:survival-tilt-def} is justified by a
uniform exponential majorant in a neighborhood of each real $s$.  It gives
\begin{equation}
 J_k'(s)=-J_{k+1}(s).
 \label{eq:J-differential-recurrence}
\end{equation}
Substituting into \eqref{eq:M-zero}--\eqref{eq:M-successor} yields the cleaner recurrence
\begin{equation}
 M_k'(s)=-M_{k+1}(s).
 \label{eq:M-differential-recurrence}
\end{equation}
Iteration gives
\begin{equation}
 M_k^{(n)}(s)=(-1)^nM_{k+n}(s).
 \label{eq:M-iterated-derivative}
\end{equation}
Thus all $M_k$ are $C^\infty$ on $\R$, and the negative generating function is completely
monotone:
\begin{equation}
 (-1)^n\frac{\dd^n}{\dd s^n}G(-s)=M_n(s)\ge0.
 \label{eq:complete-monotonicity}
\end{equation}

\subsection{Reflection and the half moments}

Reflection invariance \eqref{eq:law-reflection} gives
\begin{equation}
 \int(1-x)^n\dd\mu(x)=\int x^n\dd\mu(x).
 \label{eq:endpoint-raw-reflection}
\end{equation}
At zero tilt the analytic moment normalization is the exact rational half-moment sequence
$d_n$:
\begin{equation}
 M_n(0)=d_n.
 \label{eq:M-zero-half-moment}
\end{equation}
Combining the identifications yields the full probability interpretation.

\begin{theorem}[Probability meaning of the half moments]
\label{thm:half-moments-probability}
For every $n\in\N$,
\begin{equation}
 d_n=\E[X^n]=\E[(1-X)^n]
     =\int_{[0,1]}x^n\dd\mu(x)
     =\int_{[0,1]}(1-x)^n\dd\mu(x).
 \label{eq:half-moment-probability}
\end{equation}
Moreover
\begin{equation}
 \left.\frac{\dd^n}{\dd s^n}G(-s)\right|_{s=0}=(-1)^nd_n.
 \label{eq:generating-jet-half-moments}
\end{equation}
\end{theorem}

\begin{proof}
Use \eqref{eq:law-M-identification} at $s=0$, then
\eqref{eq:M-zero-half-moment} and \eqref{eq:endpoint-raw-reflection}.  The derivative identity
is \eqref{eq:M-iterated-derivative} with $k=0$ and $s=0$.
\end{proof}

\begin{boxedremark}
The equality $d_n=\E[X^n]$ is not merely an interpretation of a previously known rational
sequence.  It is the bridge that licenses positivity, log-convexity, moment-matrix
positivity, and Laplace comparison arguments without rebuilding them from the recurrence for
$d_n$.  The formal library keeps the bridge in a separate module precisely so the
probabilistic and saddle branches can import it independently.
\end{boxedremark}

\section{From atomic laws to half-endpoint histograms}
\label{sec:step-measure-bridge}

The main article introduces the polynomial step approximants $\Wcal_n$ and then moves quickly
to their convergence.  The formal proof has a nontrivial middle layer: the coefficients first
define an atomic probability measure; each atom is smeared across a centered cell; interval
integrals of the resulting histogram are sandwiched between masses of slightly shrunken and
enlarged intervals.

\subsection{Half-endpoint indicators are ordinary indicators almost everywhere}

For $a<b$, define
\begin{equation}
 \mathbf1_{[a,b]}^*(x)=
 \begin{cases}
  1,&a<x<b,\\
  \tfrac12,&x=a\text{ or }x=b,\\
  0,&\text{otherwise}.
 \end{cases}
 \label{eq:half-endpoint-indicator}
\end{equation}
The two endpoint values are essential for pointwise gluing of adjacent cells, but irrelevant
to Lebesgue integration.

\begin{lemma}[Almost-everywhere replacement]
\label{lem:half-indicator-ae}
For every $a,b\in\R$,
\begin{equation}
 \mathbf1_{[a,b]}^*=\mathbf1_{(a,b]}
 \quad\text{Lebesgue-almost everywhere}.
 \label{eq:half-indicator-ae}
\end{equation}
Consequently
\begin{equation}
 \int_\R\mathbf1_{[a,b]}^*(x)\dd x=\max(b-a,0),
 \label{eq:half-indicator-total-integral}
\end{equation}
and, for $c\le d$,
\begin{equation}
 \int_c^d\mathbf1_{[a,b]}^*(x)\dd x
 =\max\bigl(\min(d,b)-\max(c,a),0\bigr).
 \label{eq:half-indicator-overlap}
\end{equation}
\end{lemma}

\begin{proof}
The functions differ only at $a$ and $b$, a null set.  The first integral is the length of
$(a,b]$.  The second is the length of
$(c,d]\cap(a,b]=(\max(c,a),\min(d,b)]$, with truncation at zero when the interval is empty.
\end{proof}

\subsection{Atomic measure, cell geometry, and exact normalization}

At level $n$, let $x_{n,m}$ be the location of the $m$th polynomial atom and let
$p_{n,m}\ge0$ be its normalized mass, with
\begin{equation}
 \sum_m p_{n,m}=1.
 \label{eq:atomic-mass-one}
\end{equation}
The corresponding atomic law is
\begin{equation}
 \nu_n=\sum_m p_{n,m}\delta_{x_{n,m}}.
 \label{eq:polynomial-atomic-law}
\end{equation}
Let
\begin{equation}
 h_n=2^{-n-1}
 \label{eq:step-half-width}
\end{equation}
be the half-width of a cell and
\begin{equation}
 C_{n,m}=[x_{n,m}-h_n,x_{n,m}+h_n].
 \label{eq:step-cell}
\end{equation}
Its full width is $2h_n=2^{-n}$.  Smearing the atom uniformly over its cell produces the
histogram
\begin{equation}
 \Wcal_n(x)=2^n\sum_m p_{n,m}\mathbf1^*_{C_{n,m}}(x).
 \label{eq:histogram-from-atoms}
\end{equation}
By \eqref{eq:half-indicator-total-integral}, each term has integral $p_{n,m}$, hence
\begin{equation}
 \int_\R\Wcal_n(x)\dd x=1.
 \label{eq:histogram-mass-one}
\end{equation}
The pointwise half weights ensure that when two cells share an endpoint, the two half
contributions add to the correct plateau value.

\subsection{The overlap sandwich}

For a fixed interval $[a,b]$ with $a\le b$, let
\begin{equation}
 \ell_{n,m}(a,b)=|[a,b]\cap C_{n,m}|
 =\max\bigl(\min(b,x_{n,m}+h_n)-\max(a,x_{n,m}-h_n),0\bigr).
 \label{eq:cell-overlap-length}
\end{equation}
Then
\begin{equation}
 \int_a^b\Wcal_n(x)\dd x
 =\sum_m p_{n,m}\,2^n\ell_{n,m}(a,b).
 \label{eq:histogram-overlap-sum}
\end{equation}
The normalized overlap $2^n\ell_{n,m}$ lies in $[0,1]$.

If $x_{n,m}\in(a+h_n,b-h_n]$, then the entire cell lies inside $[a,b]$, so the normalized
overlap is $1$.  If the cell meets $[a,b]$, then its center lies in
$(a-h_n,b+h_n]$.  Therefore, term by term,
\begin{equation}
 \mathbf1_{(a+h_n,b-h_n]}(x_{n,m})
 \le2^n\ell_{n,m}(a,b)
 \le\mathbf1_{(a-h_n,b+h_n]}(x_{n,m}).
 \label{eq:overlap-indicator-sandwich}
\end{equation}
Multiplication by $p_{n,m}$ and summation gives the central bridge.

\begin{theorem}[Atomic--histogram interval sandwich]
\label{thm:atomic-histogram-sandwich}
For $a\le b$,
\begin{equation}
 \nu_n((a+h_n,b-h_n])
 \le\int_a^b\Wcal_n(x)\dd x
 \le\nu_n((a-h_n,b+h_n]).
 \label{eq:atomic-histogram-sandwich}
\end{equation}
\end{theorem}

This is the exact content of the inner- and outer-overlap lemmas in
\path{StepMeasureBridge.lean}.  Its importance is conceptual: it converts weak convergence
of measures into convergence of integrals of discontinuous histograms without claiming
pointwise convergence of densities from weak convergence alone.

\subsection{Portmanteau plus a moving-boundary squeeze}

The atomic laws $\nu_n$ coincide with the finite convolution laws and converge weakly to the
absolutely continuous law with density $\Up$.  Because that limit gives zero mass to every
singleton, the frontier of $(a,b]$ is null.  Portmanteau therefore gives
\begin{equation}
 \nu_n((a,b])\longrightarrow\int_a^b\Up(x)\dd x.
 \label{eq:portmanteau-Ioc}
\end{equation}
The intervals in \eqref{eq:atomic-histogram-sandwich} move with $n$, so one more squeeze is
needed.  Fix $\varepsilon>0$ and choose $\delta>0$ with $4\delta<\varepsilon$.  Eventually
$h_n\le\delta$, and hence
\begin{align}
 (a+\delta,b-\delta]&\subset(a+h_n,b-h_n],                 \label{eq:fixed-inner-subset}\\
 (a-h_n,b+h_n]&\subset(a-\delta,b+\delta].                 \label{eq:moving-outer-subset}
\end{align}
The bound $0\le\Up\le1$ gives
\begin{align}
 0&\le\int_a^b\Up-\int_{a+\delta}^{b-\delta}\Up\le2\delta, \label{eq:up-shrink-error}\\
 0&\le\int_{a-\delta}^{b+\delta}\Up-\int_a^b\Up\le2\delta. \label{eq:up-expand-error}
\end{align}
Apply \eqref{eq:portmanteau-Ioc} to the two fixed comparison intervals and combine with
\eqref{eq:atomic-histogram-sandwich}.

\begin{theorem}[Interval-integral convergence of the step approximants]
\label{thm:step-integral-convergence}
For every $a,b\in\R$,
\begin{equation}
 \int_a^b\Wcal_n(x)\dd x
 \longrightarrow
 \int_a^b\Up(x)\dd x.
 \label{eq:step-integral-convergence}
\end{equation}
\end{theorem}

\begin{proof}
The preceding fixed-$\delta$ sandwich proves the result for $a\le b$.  If $b<a$, reverse the
orientation of both interval integrals.
\end{proof}

\begin{warning}[The logical order cannot be shortened]
Weak convergence of $\nu_n$ does not by itself imply pointwise convergence of the
histograms $\Wcal_n$, nor even convergence of their values at a single point.  The formal
route is
\[
\begin{aligned}
 \text{weak convergence}
 &\Longrightarrow \text{null-boundary interval masses}\\
 &\Longrightarrow \text{moving-cell sandwich}
 \Longrightarrow \text{interval-integral convergence}.
\end{aligned}
\]
Other parts of the development then use shape, refinement, Fourier decay, or difference
quotients to obtain stronger function-level convergence.
\end{warning}

\part{The reusable discrete engines}

\section{Weighted paths solve triangular recurrences}
\label{sec:path-sums}

The exact inverse-power recurrence for $F(2^{-n})$ is triangular.  The main article states its
composition formula, but the Lean proof first establishes a generic solver for every
semiring-valued triangular recurrence.  That abstraction explains both the form of the
answer and the uniqueness argument.

\subsection{Compositions as increasing paths}

A composition of $n$ is an ordered list of positive integers
\begin{equation}
 c=(r_1,\ldots,r_m),
 \qquad
 r_1+\cdots+r_m=n.
 \label{eq:composition}
\end{equation}
Its partial sums
\begin{equation}
 s_0=0,
 \qquad
 s_j=r_1+\cdots+r_j
 \quad(1\le j\le m)
 \label{eq:composition-partial-sums}
\end{equation}
form a strictly increasing path
\begin{equation}
 0=s_0<s_1<\cdots<s_m=n.
 \label{eq:increasing-path}
\end{equation}
Conversely, every such path determines the positive increments $r_j=s_j-s_{j-1}$.

Let $R$ be a semiring and let
\begin{equation}
 w:\N\times\N\to R
 \label{eq:generic-edge-weight}
\end{equation}
be an edge-weight function.  Define
\begin{align}
 W_w(c)&=\prod_{j=1}^{m}w(s_{j-1},s_j),                    \label{eq:path-weight}\\
 P_w(n)&=\sum_{c\models n}W_w(c).                          \label{eq:path-sum}
\end{align}
For $n=0$, there is one empty composition and its empty product is $1$, so
\begin{equation}
 P_w(0)=1.
 \label{eq:path-sum-zero}
\end{equation}

\subsection{Last-edge decomposition}

Every nonempty path to $n$ has a unique last predecessor $k<n$.  Deleting its final edge
gives a path to $k$, and appending the edge $k\to n$ reverses the operation.  Therefore

\begin{theorem}[Last-edge decomposition]
\label{thm:path-last-edge}
For $n>0$,
\begin{equation}
 P_w(n)=\sum_{k=0}^{n-1}P_w(k)w(k,n).
 \label{eq:path-last-edge}
\end{equation}
\end{theorem}

\begin{proof}
Partition the finite set of compositions of $n$ by the sum $k$ of all blocks except the last.
For fixed $k$, deleting the final block is a bijection with compositions of $k$.  Path-weight
multiplicativity gives
\[
 W_w(c\mathbin{\|}(n-k))=W_w(c)w(k,n).
\]
Summing first over compositions of $k$ and then over $k<n$ gives
\eqref{eq:path-last-edge}.
\end{proof}

One can instead partition by the number of edges.  Let $P_{w,r}(n)$ be the contribution of
compositions with exactly $r$ blocks.  Uniformly in $n$,
\begin{equation}
 P_w(n)=\sum_{r=0}^{n}P_{w,r}(n),
 \label{eq:path-length-uniform}
\end{equation}
and for $n>0$ the zero-edge term vanishes, so the range may be written $1\le r\le n$.

\subsection{The generic triangular solver}

\begin{theorem}[Path-sum solution of a triangular recurrence]
\label{thm:triangular-path-solver}
Let $x:\N\to R$ satisfy
\begin{align}
 x_0&=b,                                                    \label{eq:triangular-initial}\\
 x_n&=\sum_{k=0}^{n-1}x_k w(k,n)\qquad(n>0).               \label{eq:triangular-recurrence}
\end{align}
Then
\begin{equation}
 x_n=bP_w(n)
 \qquad(n\ge0).
 \label{eq:triangular-path-solution}
\end{equation}
Consequently the recurrence and initial value determine $x$ uniquely.
\end{theorem}

\begin{proof}
Use strong induction on $n$.  At $n=0$, \eqref{eq:path-sum-zero} gives
$bP_w(0)=b=x_0$.  For $n>0$, insert the induction hypotheses into
\eqref{eq:triangular-recurrence}:
\[
 x_n=\sum_{k<n}bP_w(k)w(k,n)
 =b\sum_{k<n}P_w(k)w(k,n)
 =bP_w(n),
\]
where the last equality is \cref{thm:path-last-edge}.  Any two solutions therefore equal the
same path sum at each index.
\end{proof}

\begin{remark}
No subtraction, division, topology, or order is used.  The theorem works in an arbitrary
semiring.  This is why the formal development separates it from the rational Fabius
specialization.
\end{remark}

\section{The nonrecursive composition formula for \texorpdfstring{$F(2^{-n})$}{F(2 to the -n)}}
\label{sec:fabius-composition-formula}

Let $a_n$ be the normalized inverse-dyadic recurrence sequence used in the main article, so
\begin{equation}
 F(2^{-n})=2^{-\binom n2}a_n.
 \label{eq:inverse-dyadic-normalization}
\end{equation}
Its triangular recurrence has edge weights
\begin{equation}
 w(k,n)=\frac{1}{(2^n-1)(n-k+1)!}
 \qquad(0\le k<n).
 \label{eq:fabius-recurrence-edge}
\end{equation}
Because $a_0=1$, \cref{thm:triangular-path-solver} gives $a_n=P_w(n)$.

For a composition $c=(r_1,\ldots,r_m)$ of $n$, the edge from $s_{j-1}$ to $s_j$ has
$n-k$ replaced by $r_j$, hence
\begin{equation}
 w(s_{j-1},s_j)=\frac{1}{(2^{s_j}-1)(r_j+1)!}.
 \label{eq:fabius-composition-edge}
\end{equation}
Therefore:

\begin{theorem}[Nonrecursive inverse-dyadic formula]
\label{thm:composition-inverse-dyadic}
For every $n\in\N$,
\begin{equation}
 F(2^{-n})=
 2^{-\binom n2}
 \sum_{(r_1,\ldots,r_m)\models n}
 \prod_{j=1}^{m}
 \frac{1}{(2^{r_1+\cdots+r_j}-1)(r_j+1)!}.
 \label{eq:composition-inverse-dyadic}
\end{equation}
At $n=0$, the unique empty composition contributes the empty product $1$.
\end{theorem}

\begin{proof}
Apply \cref{thm:triangular-path-solver} with \eqref{eq:fabius-recurrence-edge}, translate
paths into compositions using \eqref{eq:composition-partial-sums}, and multiply by the outer
normalization in \eqref{eq:inverse-dyadic-normalization}.
\end{proof}

The formula may be grouped by the number $m$ of blocks:
\begin{equation}
 F(2^{-n})=
 2^{-\binom n2}
 \sum_{m=1}^{n}
 \sum_{\substack{r_1+\cdots+r_m=n\\r_j\ge1}}
 \prod_{j=1}^{m}
 \frac{1}{(2^{r_1+\cdots+r_j}-1)(r_j+1)!},
 \qquad n>0.
 \label{eq:composition-by-length}
\end{equation}

\subsection{Small cases as structural checks}

\begin{example}[The value at $1/2$]
For $n=1$ the sole composition is $(1)$, with weight
\[
 \frac1{(2^1-1)2!}=\frac12.
\]
The outer power is $2^0=1$, so $F(1/2)=1/2$.
\end{example}

\begin{example}[The value at $1/4$]
The two compositions of $2$ are $(2)$ and $(1,1)$:
\begin{align*}
 W(2)&=\frac1{(4-1)3!}=\frac1{18},\\
 W(1,1)&=\frac1{(2-1)2!}\frac1{(4-1)2!}=\frac1{12}.
\end{align*}
Their sum is $5/36$ and the outer factor is $2^{-1}$, giving
\begin{equation}
 F(1/4)=\frac5{72}.
 \label{eq:F-quarter-path-check}
\end{equation}
\end{example}

\begin{example}[The value at $1/8$]
The four compositions of $3$ have weights
\begin{equation*}
 \frac1{168},\qquad
 \frac1{84},\qquad
 \frac1{252},\qquad
 \frac1{168},
\end{equation*}
for $(3)$, $(1,2)$, $(2,1)$, and $(1,1,1)$ respectively.  Their sum is $1/36$; multiplying by
$2^{-\binom32}=1/8$ gives
\begin{equation}
 F(1/8)=\frac1{288}.
 \label{eq:F-eighth-path-check}
\end{equation}
\end{example}

\begin{boxedremark}
The composition formula is not a mysterious closed form guessed from the first values.  It
is the universal unrolling of a triangular recurrence.  Every block records a jump, every
partial sum records the endpoint of that jump, and the product records the chronological
edge weights.  The formal path layer makes this explanation exact and reusable.
\end{boxedremark}

\section{A finite-row Toeplitz convergence theorem}
\label{sec:toeplitz}

The complex-shift and discrete-limit formulas involve a row of coefficients depending on
$n$, not a fixed summability kernel.  The main article verifies the Fabius specialization but
omits the abstract theorem.  The needed result is a finite-row version of the Toeplitz
summability principle in a normed ring.

\subsection{Abstract statement}

Let $K$ be a normed ring.  For each $n$, let $J_n$ be a finite index set, let
$w_{n,j}\in K$, and let $q(n,j)\in\N$ be the sampled sequence index.  Assume:
\begin{align}
 \sum_{j\in J_n}w_{n,j}&=1,                                \label{eq:toeplitz-row-mass}\\
 \sum_{j\in J_n}\|w_{n,j}\|&\le C                         \label{eq:toeplitz-variation}
\end{align}
for one constant $C\ge0$, and assume that the sampled indices move uniformly into every
tail:
\begin{equation}
 \forall N\ \exists n_0\ \forall n\ge n_0\ \forall j\in J_n,
 \qquad q(n,j)\ge N.
 \label{eq:toeplitz-tail-index}
\end{equation}

\begin{theorem}[Finite-row Toeplitz convergence]
\label{thm:finite-row-toeplitz}
If $H_m\to L$ in $K$, then
\begin{equation}
 \sum_{j\in J_n}w_{n,j}H_{q(n,j)}\longrightarrow L.
 \label{eq:finite-row-toeplitz}
\end{equation}
\end{theorem}

\begin{proof}
Fix $\varepsilon>0$.  Since $H_m\to L$, choose $N$ such that
\begin{equation}
 \|H_m-L\|<\frac{\varepsilon}{C+1}
 \qquad(m\ge N).
 \label{eq:toeplitz-tail-error}
\end{equation}
Choose $n_0$ from \eqref{eq:toeplitz-tail-index}.  For $n\ge n_0$, use the row-mass identity
to subtract $L$:
\begin{align*}
 \sum_{j\in J_n}w_{n,j}H_{q(n,j)}-L
 &=\sum_{j\in J_n}w_{n,j}\bigl(H_{q(n,j)}-L\bigr).
\end{align*}
Therefore
\begin{align*}
 \left\|\sum_{j\in J_n}w_{n,j}H_{q(n,j)}-L\right\|
 &\le\sum_{j\in J_n}\|w_{n,j}\|\,\|H_{q(n,j)}-L\|\\
 &<\frac{\varepsilon}{C+1}\sum_{j\in J_n}\|w_{n,j}\|\\
 &\le\varepsilon\frac{C}{C+1}<\varepsilon.
\end{align*}
This proves convergence.
\end{proof}

\begin{remark}
The weights need not be positive, the rows need not be nested, and no entrywise limit of
$w_{n,j}$ is required.  The theorem uses exactly three facts: row mass one, uniformly bounded
total variation, and uniform migration of all sampled indices into the tail.
\end{remark}

\section{The Fabius Toeplitz rows and corrected discrete prefixes}
\label{sec:fabius-toeplitz}

\subsection{The half-base \texorpdfstring{$q$}{q}-binomial weights}

For $0\le j\le n$, define
\begin{equation}
 w_{n,j}=\frac{(-1)^j
 \qbinom{n}{j}{1/2}
 (1/2)^{\binom{j+1}{2}}}{(1/2;1/2)_n}.
 \label{eq:fabius-toeplitz-weight}
\end{equation}
The half-base finite $q$-binomial theorem evaluated at $z=1/2$ gives
\begin{equation}
 \sum_{j=0}^{n}w_{n,j}=1.
 \label{eq:fabius-toeplitz-mass}
\end{equation}
Taking absolute values removes the signs and changes the finite product from
$(1/2;1/2)_n$ to $(-1/2;1/2)_n$:
\begin{equation}
 \sum_{j=0}^{n}|w_{n,j}|
 =\frac{(-1/2;1/2)_n}{(1/2;1/2)_n}.
 \label{eq:fabius-toeplitz-exact-variation}
\end{equation}
The formal proof uses the elementary uniform lower bound
\begin{equation}
 (1/2;1/2)_n\ge\frac14
 \label{eq:half-q-lower-quarter}
\end{equation}
and the product comparison
\begin{equation}
 (-1/2;1/2)_n(1/2;1/2)_n\le1.
 \label{eq:q-product-comparison}
\end{equation}
Hence
\begin{equation}
 \sum_{j=0}^{n}|w_{n,j}|\le16.
 \label{eq:fabius-toeplitz-variation-16}
\end{equation}
The constant $16$ is deliberately coarse but uniform and algebraically simple.

The reindexing used by the discrete formula is
\begin{equation}
 q(n,j)=2n-j.
 \label{eq:fabius-toeplitz-index}
\end{equation}
Since $0\le j\le n$,
\begin{equation}
 q(n,j)=2n-j\ge n.
 \label{eq:fabius-toeplitz-index-tail}
\end{equation}
Thus every index in row $n$ lies beyond the common threshold $n$.

\begin{corollary}[Specialized Fabius row convergence]
\label{cor:fabius-toeplitz-convergence}
For every convergent sequence $H_m\to L$ in $\R$ or $\C$,
\begin{equation}
 \sum_{j=0}^{n}w_{n,j}H_{2n-j}\longrightarrow L.
 \label{eq:fabius-toeplitz-convergence}
\end{equation}
\end{corollary}

\begin{proof}
Apply \cref{thm:finite-row-toeplitz} with $C=16$ and use
\eqref{eq:fabius-toeplitz-mass}, \eqref{eq:fabius-toeplitz-variation-16}, and
\eqref{eq:fabius-toeplitz-index-tail}.
\end{proof}

\subsection{Why the inner range uses a half-cell correction}

The finite discrete approximant contains an inner prefix whose source-facing notation has
inclusive upper endpoint
\begin{equation}
 \left\lfloor2^px-\frac12\right\rfloor.
 \label{eq:inclusive-upper-bound}
\end{equation}
A direct natural-number range cannot represent the value $-1$, which should mean an empty
sum.  The safe range length is
\begin{equation}
 \ell_p(x)=\left\lfloor2^px+\frac12\right\rfloor_{\N},
 \label{eq:corrected-range-length}
\end{equation}
where the natural floor clamps negative results to zero.  For $x\ge0$,
\begin{equation}
 \ell_p(x)=\left\lfloor2^px-\frac12\right\rfloor+1,
 \label{eq:range-length-inclusive}
\end{equation}
so the natural range $0\le r<\ell_p(x)$ exactly represents the intended inclusive sum.
Moreover
\begin{equation}
 \ell_p(x)=0
 \iff 2^px<\frac12.
 \label{eq:range-length-zero-iff}
\end{equation}
In particular every prefix is empty for $x\le0$, so every finite approximant vanishes there
before taking a limit.

\begin{warning}[The half-cell is not cosmetic]
Replacing \eqref{eq:corrected-range-length} by $\lfloor2^px\rfloor$ shifts every jump and
breaks the centered nearest-integer convention.  Replacing it by a signed inclusive endpoint
inside a natural range mishandles the empty case.  The formula
$\lfloor2^px+1/2\rfloor_{\N}$ simultaneously fixes both issues and is the exact convention
used by the formal development.
\end{warning}

\subsection{Where the abstract theorem enters the discrete limit}

After the algebraic $q$-binomial reindexing, the $n$th discrete approximant takes the form
\begin{equation}
 A_n(x)=\sum_{j=0}^{n}w_{n,j}H_{2n-j}(x),
 \label{eq:discrete-limit-toeplitz-form}
\end{equation}
where, for each fixed $x$, the inner objects $H_p(x)$ converge to the desired translated
Fabius value.  The entire outer-limit argument is then
\begin{equation*}
 H_p(x)\to L(x)
 \quad+\quad
 \text{\eqref{eq:fabius-toeplitz-mass}, \eqref{eq:fabius-toeplitz-variation-16},
 \eqref{eq:fabius-toeplitz-index-tail}}
 \quad\Longrightarrow\quad
 A_n(x)\to L(x).
\end{equation*}
All Fabius-specific work lies in identifying $H_p$ and its limit; all row summation is
contained in \cref{thm:finite-row-toeplitz}.

\part{The formal and analytic machinery of the all-orders saddle expansion}

\section{Parameter-dependent Poincar\'e expansions}
\label{sec:poincare-api}

The endpoint asymptotic is not an ordinary expansion with constant coefficients.  Its
coefficients retain a bounded one-periodic dependence on the exact Lambert phase.  The
formal library therefore uses an asymptotic API in which coefficient families may depend on
the asymptotic parameter, provided that dependence remains bounded.

Let $\ell$ be a filter on a parameter space $A$, let $\sigma:A\to\mathbb K$ be a scale tending
to zero, and let $f:A\to E$ take values in a normed vector space over the normed field
$\mathbb K$.  For coefficient families $c_k:A\to E$, put
\begin{equation}
 S_N[f](a)=\sum_{k=0}^{N-1}\sigma(a)^k c_k(a).
 \label{eq:parameter-partial-sum}
\end{equation}

\begin{definition}[Full parameter-dependent expansion]
\label{def:parameter-expansion}
We write
\begin{equation}
 f(a)\sim_{\ell}\sum_{k\ge0}\sigma(a)^kc_k(a)
 \label{eq:parameter-expansion-symbol}
\end{equation}
when
\begin{enumerate}
 \item every $c_k$ is bounded along $\ell$, equivalently $c_k=O_{\ell}(1)$;
 \item for every $N\in\N$,
 \begin{equation}
  f-S_N[f]=O_{\ell}(\sigma^N).
  \label{eq:parameter-expansion-remainder}
 \end{equation}
\end{enumerate}
\end{definition}

This is \lean{HasAsymptoticExpansion}.  Boundedness of the coefficients is not redundant: if
arbitrary unbounded parameter dependence were allowed, powers of $\sigma$ could be moved
between the ``coefficient'' and the scale, destroying the meaning of order.

\subsection{Elementary algebra of full expansions}

Suppose $f$ and $g$ have expansions with the same scale.  Direct manipulation of finite
partial sums proves:
\begin{align}
 f+g&\sim\sum_{k\ge0}\sigma^k(c_k+d_k),                    \label{eq:expansion-add}\\
 f-g&\sim\sum_{k\ge0}\sigma^k(c_k-d_k),                    \label{eq:expansion-sub}\\
 \alpha f&\sim\sum_{k\ge0}\sigma^k(\alpha c_k).            \label{eq:expansion-smul}
\end{align}
The API also supports transport across eventual equality and pullback along a map tending to
the source filter.  These apparently routine lemmas are what let the final proof switch
between complex and real saddle ratios, replace exact identities after a threshold, and move
from the logarithmic coordinate to $x\downarrow0$ without reopening every remainder
estimate.

\subsection{Flat functions}

\begin{definition}[Flat relative to a scale]
\label{def:flat-scale}
A function $r$ is flat relative to $\sigma$ along $\ell$ if
\begin{equation}
 r=O_{\ell}(\sigma^N)
 \qquad\text{for every }N\in\N.
 \label{eq:flat-all-orders}
\end{equation}
\end{definition}

\begin{theorem}[Flat remainders are invisible to all coefficients]
\label{thm:flat-invisible}
A flat $r$ has the full expansion with every coefficient zero.  Consequently, if
\begin{equation}
 f\sim\sum_{k\ge0}\sigma^kc_k,
 \label{eq:f-exp-before-flat}
\end{equation}
then
\begin{equation}
 f+r\sim\sum_{k\ge0}\sigma^kc_k,
 \qquad
 f-r\sim\sum_{k\ge0}\sigma^kc_k.
 \label{eq:flat-preserves-expansion}
\end{equation}
The implications hold in both directions.
\end{theorem}

\begin{proof}
For the zero-coefficient expansion, the $N$th partial sum is identically zero, so the
remainder condition is exactly \eqref{eq:flat-all-orders}.  Add or subtract this expansion
coefficientwise from the expansion of $f$.  Reversibility follows by moving $r$ to the other
side.
\end{proof}

This theorem is the formal reason that the forward negative-Laplace tail can be restored
after the central saddle expansion without changing any algebraic-order coefficient.

\subsection{Exact logarithmic coordinates}

Define
\begin{equation}
 x(t)=2^{-t},
 \qquad
 t(x)=-\log_2x
 \quad(x>0).
 \label{eq:small-log-coordinates}
\end{equation}
These maps are exact inverses on $x>0$, and
\begin{equation}
 t\to+\infty\iff x(t)\downarrow0,
 \qquad
 x\downarrow0\iff t(x)\to+\infty.
 \label{eq:coordinate-filter-equivalence}
\end{equation}

\begin{theorem}[Full-expansion coordinate equivalence]
\label{thm:small-log-expansion-equivalence}
For functions $f,\sigma,c_k$ on the positive half-line,
\begin{equation}
 f(2^{-t})\sim_{t\to\infty}
 \sum_{k\ge0}\sigma(2^{-t})^kc_k(2^{-t})
 \label{eq:log-coordinate-expansion}
\end{equation}
if and only if
\begin{equation}
 f(x)\sim_{x\downarrow0}
 \sum_{k\ge0}\sigma(x)^kc_k(x).
 \label{eq:small-coordinate-expansion}
\end{equation}
\end{theorem}

\begin{proof}
Pull back the full expansion along either coordinate map.  The two composites are eventually
identical on their respective filters by \eqref{eq:small-log-coordinates}; boundedness and
every remainder estimate therefore transfer in both directions.
\end{proof}

\section{Recursive coefficients of a formal exponential}
\label{sec:exp-coeff}

The saddle integrand contains the exponential of a formal perturbation.  Rather than invoke
Bell polynomials or repeatedly expand a universal power series, the Lean development uses a
finite recurrence whose $n$th output depends only on the first $n$ input coefficients.

Let $R$ be a commutative $\Q$-algebra, and let
\begin{equation}
 E(u)=\sum_{j\ge1}E_j u^j
 \label{eq:formal-exponent}
\end{equation}
be a formal exponent with zero constant term.  Define $a_n\in R$ recursively by
\begin{align}
 a_0&=1,                                                    \label{eq:expcoeff-zero}\\
 a_{n+1}&=\frac1{n+1}\sum_{j=0}^{n}(j+1)E_{j+1}a_{n-j}.
                                                               \label{eq:expcoeff-recurrence}
\end{align}

\begin{theorem}[Correctness and uniqueness of the exponential recurrence]
\label{thm:expcoeff-correctness}
The power series
\begin{equation}
 A(u)=\sum_{n\ge0}a_nu^n
 \label{eq:formal-mass-series}
\end{equation}
satisfies
\begin{equation}
 A'(u)=E'(u)A(u),
 \qquad
 A(0)=1.
 \label{eq:formal-exp-ode}
\end{equation}
It is the unique formal power series with these two properties, and therefore
\begin{equation}
 A(u)=\exp(E(u)).
 \label{eq:formal-exp-identification}
\end{equation}
\end{theorem}

\begin{proof}
The coefficient of $u^n$ in $A'$ is $(n+1)a_{n+1}$.  The coefficient of $u^n$ in $E'A$ is
\[
 \sum_{j=0}^{n}(j+1)E_{j+1}a_{n-j},
\]
so \eqref{eq:formal-exp-ode} is coefficientwise equivalent to
\eqref{eq:expcoeff-recurrence}.  Conversely the differential equation and initial value
determine $a_0$, then $a_1$, then $a_2$, and so on, proving uniqueness.  The universal
formal exponential satisfies the same equation and initial value.
\end{proof}

The first coefficients are
\begin{align}
 a_1={}&E_1,                                                  \label{eq:a1}\\
 a_2={}&E_2+\frac12E_1^2,                                    \label{eq:a2}\\
 a_3={}&E_3+E_1E_2+\frac16E_1^3,                             \label{eq:a3}\\
 a_4={}&E_4+E_1E_3+\frac12E_2^2
       +\frac12E_1^2E_2+\frac1{24}E_1^4.                    \label{eq:a4}
\end{align}
Equation \eqref{eq:a4} is the exact finite identity used for the second saddle correction.

\subsection{Structural laws}

The recurrence is preferable to a black-box exponential because it exposes several laws
needed downstream.

\begin{proposition}[Finite dependence]
\label{prop:expcoeff-finite-dependence}
The coefficient $a_n$ depends only on $E_1,\ldots,E_n$.  If two exponent families agree
through order $n$, their $n$th exponential coefficients agree.
\end{proposition}

\begin{proposition}[Naturality]
\label{prop:expcoeff-naturality}
For every morphism of commutative $\Q$-algebras $\varphi:R\to S$,
\begin{equation}
 \varphi(a_n(E))=a_n(\varphi\circ E).
 \label{eq:expcoeff-map}
\end{equation}
In particular, evaluation of polynomial- or function-valued coefficients commutes with the
recurrence.
\end{proposition}

\begin{proposition}[Rescaling, sums, and parity]
\label{prop:expcoeff-structure}
Let $c\in R$ be central.
\begin{enumerate}
 \item If $\widetilde E_j=c^jE_j$, then
 \begin{equation}
  a_n(\widetilde E)=c^na_n(E).
  \label{eq:expcoeff-rescale}
 \end{equation}
 \item If $E$ and $F$ are two exponent families, then
 \begin{equation}
  a_n(E+F)=\sum_{r+s=n}a_r(E)a_s(F).
  \label{eq:expcoeff-add-convolution}
 \end{equation}
 \item If $E_j(-v)=(-1)^jE_j(v)$, then
 \begin{equation}
  a_n(-v)=(-1)^na_n(v).
  \label{eq:expcoeff-parity}
 \end{equation}
\end{enumerate}
\end{proposition}

\begin{proofidea}
Naturality is induction through finite sums and rational scalar multiplication.
Rescaling follows from the recurrence because every summand has total degree
$(j+1)+(n-j)=n+1$.  The addition law is the coefficient form of
$\exp(E+F)=\exp(E)\exp(F)$.  Parity is the specialization $c=-1$ of rescaling after
pointwise reflection in $v$.
\end{proofidea}

\section{Recursive coefficients of a formal logarithm}
\label{sec:log-coeff}

Let
\begin{equation}
 A(u)=1+\sum_{n\ge1}a_nu^n.
 \label{eq:unit-mass-series}
\end{equation}
We seek
\begin{equation}
 B(u)=\log A(u)=\sum_{n\ge1}b_nu^n.
 \label{eq:formal-log-series}
\end{equation}
Instead of expanding the universal logarithm, use
\begin{equation}
 A(u)B'(u)=A'(u).
 \label{eq:formal-log-ode}
\end{equation}
Solving the coefficient of $u^n$ for $b_{n+1}$ gives
\begin{align}
 b_0&=0,                                                     \label{eq:logcoeff-zero}\\
 b_{n+1}&=a_{n+1}-\frac1{n+1}
 \sum_{j=0}^{n-1}(n-j)b_{n-j}a_{j+1}.                      \label{eq:logcoeff-recurrence}
\end{align}
The empty sum at $n=0$ gives $b_1=a_1$.

\begin{theorem}[Correctness and uniqueness of the logarithm recurrence]
\label{thm:logcoeff-correctness}
If $a_0=1$, the series generated by \eqref{eq:logcoeff-recurrence} is the unique series with
zero constant term satisfying \eqref{eq:formal-log-ode}; hence it is the formal logarithm of
$A$.
\end{theorem}

\begin{proof}
The coefficient of $u^n$ in $AB'$ contains the term $(n+1)b_{n+1}$ from $a_0b'$, plus terms
involving only $b_1,\ldots,b_n$.  Equating it with $(n+1)a_{n+1}$ and dividing by $n+1$
gives \eqref{eq:logcoeff-recurrence}.  This triangular relation proves uniqueness.  The
formal logarithm satisfies $B(0)=0$ and $B'A=A'$.
\end{proof}

The first coefficients are
\begin{align}
 b_1={}&a_1,                                                  \label{eq:b1}\\
 b_2={}&a_2-\frac12a_1^2,                                   \label{eq:b2}\\
 b_3={}&a_3-a_1a_2+\frac13a_1^3,                            \label{eq:b3}\\
 b_4={}&a_4-a_1a_3-\frac12a_2^2+a_1^2a_2-\frac14a_1^4.     \label{eq:b4}
\end{align}
As with the exponential recurrence, $b_n$ depends only on $a_1,\ldots,a_n$, commutes with
$\Q$-algebra morphisms and evaluation, obeys the rescaling law
\begin{equation}
 a_j\mapsto c^ja_j\quad\Longrightarrow\quad b_n\mapsto c^nb_n,
 \label{eq:logcoeff-rescale}
\end{equation}
and preserves parity.

\begin{remark}[Why $a_0$ is a hypothesis, not an input]
The recursive definition \eqref{eq:logcoeff-recurrence} never reads $a_0$; it is meaningful
for every sequence.  Correctness as a logarithm, however, requires $a_0=1$.  Separating the
recurrence from its normalization makes finite coefficient calculations easier while keeping
the theorem honest.
\end{remark}

\subsection{From mass expansions to logarithmic expansions}

Suppose a positive real function has a full expansion
\begin{equation}
 R(a)\sim1+\sum_{n\ge1}\sigma(a)^na_n(a),
 \qquad \sigma(a)\to0,
 \label{eq:mass-expansion-before-log}
\end{equation}
with bounded coefficient families.  Finite polynomial remainders, discussed in the next
section, show that for each $N$,
\begin{equation}
 \log R(a)=\sum_{n=1}^{N-1}\sigma(a)^nb_n(a)+O(\sigma(a)^N),
 \label{eq:log-asymptotic-transfer}
\end{equation}
where $b_n$ is generated by \eqref{eq:logcoeff-recurrence}.  The positivity hypothesis fixes
the real branch and is eventually automatic when $R\to1$.

\section{Exact finite remainders from formal coefficient identities}
\label{sec:finite-remainder}

Formal power-series equality says that two expressions have the same coefficients below a
chosen order.  An asymptotic proof needs an evaluated identity with an explicit factor of the
first omitted power.  The module \path{SaddleExpansionFiniteRemainder.lean} provides exactly
that bridge without invoking an infinite series.

Fix $L\ge1$ and truncate the exponent:
\begin{equation}
 E_{<L}(X)=\sum_{k=1}^{L-1}E_kX^k.
 \label{eq:exponent-trunc-polynomial}
\end{equation}
Truncate the scalar exponential polynomial at the same order and substitute:
\begin{equation}
 T_L(X)=\sum_{q=0}^{L-1}\frac{E_{<L}(X)^q}{q!}.
 \label{eq:finite-exp-substitution}
\end{equation}
Let
\begin{equation}
 A_{<L}(X)=\sum_{k=0}^{L-1}a_kX^k
 \label{eq:expcoeff-trunc-polynomial}
\end{equation}
be the polynomial formed from the recursive exponential coefficients.

\begin{theorem}[Exact divisibility of the finite defect]
\label{thm:finite-defect-divisibility}
If $E_0=0$, then
\begin{equation}
 X^L\mid\bigl(T_L(X)-A_{<L}(X)\bigr).
 \label{eq:defect-divisibility}
\end{equation}
Equivalently, there is a canonical polynomial $Q_L(X)$ such that
\begin{equation}
 T_L(X)=A_{<L}(X)+X^LQ_L(X).
 \label{eq:exact-finite-remainder-polynomial}
\end{equation}
\end{theorem}

\begin{proof}
The coefficient of $X^k$ in $T_L$ agrees with that of the full formal substitution
$\exp(E(X))$ for every $k<L$: terms of scalar exponential degree $q\ge L$ cannot contribute
below degree $L$ because $E$ has zero constant term, and truncating $E$ itself does not alter
coefficients below $L$.  By \cref{thm:expcoeff-correctness}, that common coefficient is
$a_k$.  Hence the first $L$ coefficients of the defect vanish, which is equivalent to
divisibility by $X^L$.  The formal construction takes $Q_L$ to be the result of applying
polynomial division by $X$ exactly $L$ times.
\end{proof}

Evaluating at $x\in R$ gives the identity actually used in estimates:
\begin{equation}
 \sum_{q=0}^{L-1}\frac1{q!}
 \left(\sum_{k=1}^{L-1}E_kx^k\right)^q
 =\sum_{k=0}^{L-1}a_kx^k+x^LQ_L(x).
 \label{eq:finite-remainder-evaluated}
\end{equation}

\begin{boxedremark}
The quotient $Q_L$ is a polynomial in $x$ whose coefficients are algebraic expressions in
the finitely many parameters $E_1,\ldots,E_{L-1}$.  If those parameter families are uniformly
bounded, then $Q_L$ is uniformly bounded on a fixed compact $x$-window.  Thus the exact
factorization \eqref{eq:finite-remainder-evaluated} turns formal coefficient agreement into a
uniform $O(x^L)$ estimate with no appeal to convergence of an infinite formal series.
\end{boxedremark}

\section{Gaussian contraction: integration becomes polynomial algebra}
\label{sec:gaussian-contraction}

After the saddle variable is scaled by $v=\sqrt{\lambda}\,\theta$, the leading kernel is the
standard Gaussian
\begin{equation}
 \gamma(v)=\e^{-v^2/2}.
 \label{eq:standard-gaussian-kernel}
\end{equation}
Every finite reference term is $\gamma(v)$ times a complex polynomial.  The generic
contraction module turns its whole-line integral into an algebraic linear functional.

\subsection{Gaussian moments}

Define the unnormalized moments
\begin{equation}
 M_n^{\mathrm G}=\int_{\R}\e^{-v^2/2}v^n\dd v.
 \label{eq:real-gaussian-moment}
\end{equation}
Every moment is absolutely integrable.  Symmetry gives
\begin{equation}
 M_{2j+1}^{\mathrm G}=0.
 \label{eq:gaussian-odd-zero}
\end{equation}
Integration by parts, with $u=v^{n+1}$ and
$\dd w=-v\e^{-v^2/2}\dd v$, gives
\begin{equation}
 M_{n+2}^{\mathrm G}=(n+1)M_n^{\mathrm G}.
 \label{eq:gaussian-moment-recurrence}
\end{equation}
Since
\begin{equation}
 M_0^{\mathrm G}=\sqrt{2\pi},
 \label{eq:gaussian-moment-zero}
\end{equation}
the normalized moments
\begin{equation}
 \mu_n=\frac{M_n^{\mathrm G}}{\sqrt{2\pi}}
 \label{eq:normalized-gaussian-moment}
\end{equation}
satisfy
\begin{equation}
 \mu_0=1,
 \quad
 \mu_1=0,
 \quad
 \mu_{n+2}=(n+1)\mu_n.
 \label{eq:normalized-gaussian-recurrence}
\end{equation}
Thus
\begin{equation}
 \mu_{2j}=(2j-1)!!,
 \qquad
 \mu_{2j+1}=0.
 \label{eq:normalized-gaussian-closed}
\end{equation}
The first even values needed below are
\begin{equation}
 \mu_4=3,
 \quad
 \mu_6=15,
 \quad
 \mu_8=105,
 \quad
 \mu_{10}=945,
 \quad
 \mu_{12}=10395.
 \label{eq:gaussian-moments-second-correction}
\end{equation}

\subsection{The contraction functional}

For $p(X)=\sum_kp_kX^k\in\C[X]$, define
\begin{equation}
 \Gcal[p]=\sum_kp_k\mu_k.
 \label{eq:gaussian-contraction-def}
\end{equation}
This is a $\C$-linear map $\C[X]\to\C$.

\begin{theorem}[Gaussian polynomial contraction]
\label{thm:gaussian-contraction}
For every complex polynomial $p$,
\begin{equation}
 \int_{\R}\e^{-v^2/2}p(v)\dd v
 =\sqrt{2\pi}\,\Gcal[p].
 \label{eq:gaussian-contraction-integral}
\end{equation}
In particular
\begin{equation}
 \Gcal[X^{2j}]=(2j-1)!!,
 \qquad
 \Gcal[X^{2j+1}]=0.
 \label{eq:gaussian-contraction-monomials}
\end{equation}
\end{theorem}

\begin{proof}
Expand $p$ as a finite sum of monomials, interchange that finite sum with the integral, and
use \eqref{eq:normalized-gaussian-moment}.  Integrability follows coefficientwise from the
absolute Gaussian moments.
\end{proof}

The same coefficientwise expansion yields the norm estimate
\begin{equation}
 \left|\Gcal[p]\right|
 \le\sum_k|p_k|\mu_{2\lceil k/2\rceil},
 \label{eq:gaussian-contraction-norm-bound}
\end{equation}
with odd coefficients contributing zero in the exact functional.  In the formal proof a
slightly more uniform support sum is used so it composes smoothly with polynomial families.

\begin{remark}[Parity is an algebraic annihilator]
The vanishing of odd Gaussian moments is not merely a simplification after integration.  It
is the mechanism that removes every odd power of the small parameter $\varepsilon=
\lambda^{-1/2}$ from the mass expansion.  The exponent coefficient of order $m$ has parity
$m$ in $v$, the formal exponential recurrence preserves total parity, and
\eqref{eq:gaussian-contraction-monomials} kills every odd result.  The final expansion is
therefore in integer powers of $\lambda^{-1}$.
\end{remark}

\section{Uniform Gaussian tails for polynomial references}
\label{sec:gaussian-tails}

Whole-line contraction is useful only after the central-window integral has been replaced by
a whole-line integral with a negligible error.  The omitted Gaussian-tail modules provide a
single coefficient weight that controls every polynomial reference.

\subsection{A factorial coefficient weight}

For $p(X)=\sum_kp_kX^k$, define
\begin{equation}
 \mathfrak W(p)=\sum_{k\in\supp p}|p_k|\,8k!.
 \label{eq:gaussian-tail-weight}
\end{equation}
It is finite and nonnegative.

\begin{lemma}[Monomial Gaussian tail]
\label{lem:monomial-gaussian-tail}
For $A\ge4$ and $k\in\N$,
\begin{equation}
 \int_{|v|>A}\e^{-v^2/2}|v|^k\dd v
 \le 8k!\,\frac{\e^{-A^2/4}}{A}.
 \label{eq:monomial-gaussian-tail}
\end{equation}
\end{lemma}

\begin{proofidea}
Split the kernel as
$\e^{-v^2/2}=\e^{-v^2/4}\e^{-v^2/4}$.  The factor
$|v|^k\e^{-v^2/4}$ is controlled uniformly by a factorial bound, while the remaining
Gaussian tail is estimated by integrating the derivative of $\e^{-v^2/4}$ and using
$|v|\ge A$.  The constants are chosen so the result holds simultaneously for every
$k\ge0$ once $A\ge4$; no optimization is needed in the saddle application.
\end{proofidea}

\begin{theorem}[Polynomial Gaussian tail]
\label{thm:polynomial-gaussian-tail}
For every $p\in\C[X]$ and $A\ge4$,
\begin{equation}
 \int_{|v|>A}\left|\e^{-v^2/2}p(v)\right|\dd v
 \le \mathfrak W(p)\frac{\e^{-A^2/4}}{A}.
 \label{eq:polynomial-gaussian-tail}
\end{equation}
\end{theorem}

\begin{proof}
Use
$|p(v)|\le\sum_{k\in\supp p}|p_k||v|^k$, integrate the finite sum, and apply
\cref{lem:monomial-gaussian-tail} coefficientwise.
\end{proof}

\subsection{The all-orders central radius}

For expansion order $N$ and scale parameter $b\ge1$, set
\begin{equation}
 A_N(b)=\sqrt{32(N+1)\log b}.
 \label{eq:order-central-radius}
\end{equation}
Then
\begin{equation}
 \e^{-A_N(b)^2/4}=b^{-8(N+1)}.
 \label{eq:order-radius-exponential}
\end{equation}
The exponent $8(N+1)$ is much stronger than the requested algebraic order $N+1$, leaving
ample room for the factor $1/A_N(b)$ and for bounded coefficient weights.

\begin{theorem}[All-orders polynomial tail]
\label{thm:all-orders-polynomial-tail}
Let $b(a)\to+\infty$ along a filter and let $p_a\in\C[X]$ satisfy
\begin{equation}
 \mathfrak W(p_a)=O(1).
 \label{eq:bounded-polynomial-tail-weight}
\end{equation}
Then for every fixed $N$,
\begin{equation}
 \int_{|v|>A_N(b(a))}
 \left|\e^{-v^2/2}p_a(v)\right|\dd v
 =O\bigl(b(a)^{-(N+1)}\bigr).
 \label{eq:all-orders-polynomial-tail}
\end{equation}
\end{theorem}

\begin{proof}
Eventually $A_N(b(a))\ge4$ and $b(a)\ge1$.  Apply
\cref{thm:polynomial-gaussian-tail}, substitute
\eqref{eq:order-radius-exponential}, and use
\[
 b^{-8(N+1)}A_N(b)^{-1}\le b^{-(N+1)}.
\]
Multiplication by the bounded family \eqref{eq:bounded-polynomial-tail-weight} preserves the
big-$O$ rate.
\end{proof}

\begin{remark}
The theorem is formulated for a moving polynomial family.  This is essential: the saddle
reference polynomial depends periodically on the exact phase.  Periodicity and smoothness
make its finitely many coefficient functions bounded, hence its factorial weight is bounded.
\end{remark}

\section{Closed saddle jets and shifted Stirling numbers}
\label{sec:saddle-jets}

The main article states several low-order jets.  The new closed-form modules identify the
entire jet sequence and explain the combinatorial coefficients.

Let $L=\log2$ and let $\Psi$ be the smooth one-periodic fluctuation in the exact negative
Laplace exponent.  The bounded exponent jets $Z_n(t)$ satisfy
\begin{align}
 Z_0(t)&=-\frac12+\frac{\Psi'(t)}{L},                         \label{eq:Z0-recurrence}\\
 Z_{n+1}(t)&=\frac{Z_n'(t)}{L}-(n+1)Z_n(t)
              +\frac{(-1)^{n+1}n!}{L}.                      \label{eq:Z-recurrence}
\end{align}

Define integers $c(n,m)$ by
\begin{equation}
 \prod_{k=1}^{n}(z-k)=\sum_{m=0}^{n}c(n,m)z^m.
 \label{eq:falling-shift-coeff}
\end{equation}
With $H_n=1+1/2+\cdots+1/n$ and $H_0=0$, the recurrence solves explicitly.

\begin{theorem}[Closed form for every bounded exponent jet]
\label{thm:closed-saddle-jet}
For every $n\ge0$,
\begin{equation}
 Z_n(t)=(-1)^nn!\left(\frac{H_n}{L}-\frac12\right)
 +\sum_{m=0}^{n}c(n,m)
   \frac{\Psi^{(m+1)}(t)}{L^{m+1}}.
 \label{eq:closed-saddle-jet}
\end{equation}
Moreover
\begin{equation}
 c(n,m)=s(n+1,m+1),
 \label{eq:shifted-stirling}
\end{equation}
where $s(\cdot,\cdot)$ denotes the signed Stirling number of the first kind.
\end{theorem}

\begin{proof}
Let $D=L^{-1}\dd/\dd t$.  The derivative-dependent part of
\eqref{eq:Z-recurrence} is generated by the operator product
\begin{equation}
 (D-n)(D-(n-1))\cdots(D-1)\frac{\Psi'}L.
 \label{eq:jet-operator-product}
\end{equation}
Expanding the product gives the sum in \eqref{eq:closed-saddle-jet}.  The constant part
$C_n$ satisfies
\begin{equation}
 C_{n+1}=-(n+1)C_n+\frac{(-1)^{n+1}n!}{L},
 \quad C_0=-\frac12,
 \label{eq:jet-constant-recurrence}
\end{equation}
whose solution is $C_n=(-1)^nn!(H_n/L-1/2)$, using
$H_{n+1}=H_n+1/(n+1)$.  Finally, the coefficients of
$\prod_{k=1}^{n}(z-k)$ satisfy
\begin{equation}
 c(n+1,m)=c(n,m-1)-(n+1)c(n,m),
 \label{eq:c-stirling-recurrence}
\end{equation}
which is exactly the shifted signed-Stirling recurrence; the initial rows agree.
\end{proof}

The first four jets are
\begin{align}
 Z_0={}&-\frac12+\frac{\Psi'}L,                              \label{eq:Z0-explicit}\\
 Z_1={}&\frac12-\frac1L-\frac{\Psi'}L+\frac{\Psi''}{L^2},  \label{eq:Z1-explicit}\\
 Z_2={}&-1+\frac3L+\frac{2\Psi'}L-\frac{3\Psi''}{L^2}
             +\frac{\Psi^{(3)}}{L^3},                       \label{eq:Z2-explicit}\\
 Z_3={}&3-\frac{11}{L}-\frac{6\Psi'}L+\frac{11\Psi''}{L^2}
             -\frac{6\Psi^{(3)}}{L^3}+\frac{\Psi^{(4)}}{L^4}.
                                                                    \label{eq:Z3-explicit}
\end{align}
All derivatives are evaluated at the common argument $t$.  These formulas make boundedness,
one-periodicity, and $C^\infty$ regularity immediate.

\subsection{The exponent coefficient at arbitrary order}

Let $\varepsilon=\lambda^{-1/2}$ and $v$ be the Gaussian variable.  The coefficient of
$\varepsilon^m$ in the centered saddle exponent is, for $m\ge1$,
\begin{equation}
 E_m(t,v)=\ii^m\left(
 \frac{Z_{m-1}(t)}{m!}v^m
 +\frac{(-1)^{m+1}}{m+2}v^{m+2}
\right).
 \label{eq:closed-exponent-coefficient}
\end{equation}
The second monomial is universal: it is independent of $\Psi$ and of every jet.  It arises
when the jet slope $(-1)^{m}(m+1)!$, the factorial $(m+2)!$, and the extra factor
$\ii^{m+2}=-\ii^m$ cancel.

By \eqref{eq:closed-saddle-jet}, each $E_m$ is therefore completely explicit in terms of
$L^{-1}$ and $\Psi',\ldots,\Psi^{(m)}$.  It has the parity law
\begin{equation}
 E_m(t,-v)=(-1)^mE_m(t,v).
 \label{eq:exponent-coefficient-parity}
\end{equation}

\section{From exponent coefficients to periodic mass and log coefficients}
\label{sec:saddle-coeff-engine}

Let $a_n(t,v)$ be the recursive exponential coefficient
\begin{equation}
 \exp\!\left(\sum_{m\ge1}\varepsilon^mE_m(t,v)\right)
 \sim\sum_{n\ge0}\varepsilon^na_n(t,v).
 \label{eq:saddle-exp-formal}
\end{equation}
Parity preservation gives
\begin{equation}
 a_n(t,-v)=(-1)^na_n(t,v).
 \label{eq:saddle-a-parity}
\end{equation}
Define the mass coefficients by Gaussian contraction:
\begin{equation}
 B_j(t)=\Gcal[a_{2j}(t,\cdot)].
 \label{eq:saddle-mass-coefficient}
\end{equation}
Odd orders disappear.  The normalized saddle mass therefore has the formal pattern
\begin{equation}
 1+\frac{B_1(t)}{\lambda}
   +\frac{B_2(t)}{\lambda^2}+\cdots.
 \label{eq:saddle-mass-pattern}
\end{equation}
Apply the logarithm recurrence to $1+B_1u+B_2u^2+\cdots$ and define
\begin{equation}
 A_j(t)=b_j(B_0,B_1,\ldots),
 \qquad B_0=1.
 \label{eq:saddle-log-coefficient}
\end{equation}
Then
\begin{equation}
 A_0=0,
 \qquad
 A_1=B_1,
 \qquad
 A_2=B_2-\frac12B_1^2.
 \label{eq:first-log-mass-relations}
\end{equation}
Every $A_j$ and $B_j$ is bounded, one-periodic, and $C^\infty$.

\subsection{The highest derivative has a universal coefficient}

The jet $Z_{2j-1}$ first appears at mass order $j$.  The formalized leading-jet theorem says
that its coefficient in $B_j$ is
\begin{equation}
 \frac{(-1)^j}{2^jj!}.
 \label{eq:top-jet-mass-coefficient}
\end{equation}
The reason is rigid: among all monomials produced by the exponential recurrence, exactly one
can contain the newest jet, namely the linear occurrence of the exponent coefficient
$E_{2j}$.  Its jet monomial is
\begin{equation}
 \ii^{2j}\frac{Z_{2j-1}}{(2j)!}v^{2j}
 =(-1)^j\frac{Z_{2j-1}}{(2j)!}v^{2j}.
 \label{eq:top-jet-monomial}
\end{equation}
Gaussian contraction multiplies by $(2j-1)!!=(2j)!/(2^jj!)$, giving
\eqref{eq:top-jet-mass-coefficient}.

Because lower $B_k$ involve only $Z_0,\ldots,Z_{2k-1}$, every nonlinear term in the logarithm
recurrence for $A_j$ uses jets strictly below $Z_{2j-1}$.  Thus the newest jet survives with
the same coefficient.  Combining with the leading term
$\Psi^{(2j)}/L^{2j}$ in \eqref{eq:closed-saddle-jet} yields the ordinary-analysis
corollary
\begin{equation}
 [\Psi^{(2j)}]A_j
 =\frac{(-1)^j}{2^jj!L^{2j}}.
 \label{eq:top-periodic-derivative-coefficient}
\end{equation}
The mass version is a named Lean theorem; the final triangular transfer to the logarithmic
coefficient is presently a short prose corollary rather than a separate declaration.

\section{The second periodic correction, line by line}
\label{sec:second-correction}

The main article records the first correction and announces the general coefficient engine.
The current development now computes the second correction completely.  This is the first
order at which every layer of the machinery is visible at once: four exponent polynomials,
the order-four exponential coefficient, five Gaussian moments, and the mass-to-logarithm
conversion.

Fix the phase $t$ and abbreviate
\begin{equation}
 z_j=Z_j(t).
 \label{eq:zj-abbreviation}
\end{equation}
From \eqref{eq:closed-exponent-coefficient}, the first four exponent polynomials are
\begin{align}
 E_1={}&\ii\left(z_0X+\frac{X^3}{3}\right),                 \label{eq:E1-explicit}\\
 E_2={}&-\frac{z_1}{2}X^2+\frac14X^4,                       \label{eq:E2-explicit}\\
 E_3={}&\ii\left(-\frac{z_2}{6}X^3-\frac15X^5\right),      \label{eq:E3-explicit}\\
 E_4={}&\frac{z_3}{24}X^4-\frac16X^6.                      \label{eq:E4-explicit}
\end{align}
At order four in $\varepsilon$, \eqref{eq:a4} gives
\begin{equation}
 a_4=E_4+E_3E_1+\frac12E_2^2
     +\frac12E_2E_1^2+\frac1{24}E_1^4.
 \label{eq:a4-saddle}
\end{equation}
Every occurrence of $\ii$ cancels because the odd exponent polynomials carry one common
square root of $-1$.  Expanding and collecting even monomials gives
\begin{equation}
 a_4=c_4X^4+c_6X^6+c_8X^8+c_{10}X^{10}+c_{12}X^{12},
 \label{eq:a4-even-polynomial}
\end{equation}
where
\begin{align}
 c_4={}&\frac{z_0^4}{24}+\frac{z_0^2z_1}{4}
       +\frac{z_0z_2}{6}+\frac{z_1^2}{8}+\frac{z_3}{24},     \label{eq:c4}\\
 c_6={}&\frac{z_0^3}{18}-\frac{z_0^2}{8}+\frac{z_0z_1}{6}
       +\frac{z_0}{5}-\frac{z_1}{8}+\frac{z_2}{18}-\frac16, \label{eq:c6}\\
 c_8={}&\frac{z_0^2}{36}-\frac{z_0}{12}+\frac{z_1}{36}
       +\frac{47}{480},                                      \label{eq:c8}\\
 c_{10}={}&\frac{z_0}{162}-\frac1{72},                       \label{eq:c10}\\
 c_{12}={}&\frac1{1944}.                                     \label{eq:c12}
\end{align}
The five rational constants in these formulas are the visible trace of the universal
jet-free monomials in \eqref{eq:closed-exponent-coefficient}.

\subsection{Gaussian contraction}

Using \eqref{eq:gaussian-moments-second-correction}, the second mass coefficient is
\begin{equation}
 B_2=3c_4+15c_6+105c_8+945c_{10}+10395c_{12}.
 \label{eq:B2-contraction}
\end{equation}
A direct simplification gives
\begin{align}
 B_2={}&\frac{z_0^4}{8}+\frac{5z_0^3}{6}
 +\frac{3z_0^2z_1}{4}+\frac{25z_0^2}{24}
 +\frac{5z_0z_1}{2}+\frac{z_0z_2}{2}+\frac{z_0}{12}
 \notag\\
 &+\frac{3z_1^2}{8}+\frac{25z_1}{24}
 +\frac{5z_2}{6}+\frac{z_3}{8}+\frac1{288}.                \label{eq:B2-explicit}
\end{align}
The first mass coefficient, already logarithmic at first order, is
\begin{equation}
 B_1=A_1=-\frac{z_0^2}{2}-z_0-\frac{z_1}{2}-\frac1{12}.
 \label{eq:A1-jets}
\end{equation}
Substitution of \eqref{eq:Z0-explicit}--\eqref{eq:Z1-explicit} reduces this to the compact
formula from the main article,
\begin{equation}
 A_1(t)=\frac1{24}+\frac1{2L}
 -\frac{\Psi''(t)+\Psi'(t)^2}{2L^2}.
 \label{eq:A1-Psi}
\end{equation}

\subsection{Taking the logarithm}

By \eqref{eq:first-log-mass-relations},
\begin{equation}
 A_2=B_2-\frac12B_1^2.
 \label{eq:A2-from-B}
\end{equation}
The quartic and many lower-degree terms cancel.  The result is

\begin{theorem}[Closed second periodic correction]
\label{thm:closed-second-correction}
With $z_j=Z_j(t)$,
\begin{align}
 A_2(t)=\frac1{24}\bigl(&8z_0^3+12z_0^2z_1+12z_0^2
 +48z_0z_1+12z_0z_2\notag\\
 &+6z_1^2+24z_1+20z_2+3z_3\bigr).                          \label{eq:A2-jets}
\end{align}
It is a bounded smooth one-periodic function.
\end{theorem}

\begin{proof}
Insert \eqref{eq:B2-explicit} and \eqref{eq:A1-jets} into
\eqref{eq:A2-from-B}, expand the square, and collect.  Periodicity and smoothness follow from
the same properties of $Z_0,\ldots,Z_3$; boundedness follows from periodic continuity.
\end{proof}

Substitution of \eqref{eq:Z0-explicit}--\eqref{eq:Z3-explicit} makes
$A_2$ an explicit differential polynomial in
\begin{equation}
 \Psi',\ \Psi'',\ \Psi^{(3)},\ \Psi^{(4)}
 \label{eq:A2-derivatives}
\end{equation}
with coefficients in $\Q[L^{-1}]$.  The jet form \eqref{eq:A2-jets} is shorter, more stable
under formal manipulation, and displays the triangular structure that generalizes to all
orders.

\begin{boxedremark}
At order two the mass coefficient \eqref{eq:B2-explicit} contains a quartic term
$z_0^4/8$.  It disappears from $A_2$ because the logarithm removes the disconnected square
$B_1^2/2$.  This is the same connected-versus-disconnected cancellation familiar from
cumulants.  The formal logarithm recurrence performs that cancellation over an arbitrary
commutative $\Q$-algebra, before any analytic interpretation is imposed.
\end{boxedremark}

\section{Assembly of the full endpoint expansion}
\label{sec:full-assembly}

We now put the separate engines into the order in which the formal proof uses them.  This
section does not rederive the negative-Laplace product or the Bromwich formula from the main
article; it explains every transition from that exact representation to the all-orders
statement.

\subsection{Exact Lambert phase and exact normalized mass}

Let $\lambda=\lambda(x)$ be the large positive solution of
\begin{equation}
 \lambda2^{-\lambda}=x,
 \qquad x\downarrow0,
 \label{eq:exact-lambert-phase}
\end{equation}
so
\begin{equation}
 r=2^\lambda=\frac{\lambda}{x}.
 \label{eq:saddle-radius-r}
\end{equation}
Write $\Lambda(r)=\log G(-r)$ for the negative log-transform.  Define the dimensionless
normalized saddle mass
\begin{equation}
 \mathcal S(x)=
 \frac{F(x)\sqrt{2\pi\lambda}}
      {\exp(rx+\Lambda(r))}.
 \label{eq:normalized-saddle-mass}
\end{equation}
It is positive and satisfies the exact identity
\begin{equation}
 \log F(x)=rx+\Lambda(r)-\frac12\log(2\pi\lambda)+\log\mathcal S(x).
 \label{eq:exact-log-saddle-supp}
\end{equation}
No asymptotic estimate has entered yet.

Let
\begin{equation}
 \mathcal M(x)=
 -\frac L2\lambda^2+
 \left(1+\frac L2\right)\lambda-\frac12\log\lambda
 +C_F+\Psi(\lambda).
 \label{eq:sharp-main-supp}
\end{equation}
The exact negative-Laplace decomposition gives
\begin{equation}
 rx+\Lambda(r)-\frac12\log(2\pi\lambda)
 =\mathcal M(x)+E(r),
 \label{eq:exact-main-tail-supp}
\end{equation}
where the forward tail satisfies
\begin{equation}
 0\le E(2^\lambda)\le4\exp(-2^\lambda).
 \label{eq:forward-tail-flat-supp}
\end{equation}
Therefore
\begin{equation}
 \log F(x)-\mathcal M(x)=\log\mathcal S(x)+E(2^\lambda).
 \label{eq:exact-error-separation-supp}
\end{equation}
The final term is flat relative to $\lambda^{-1}$ by elementary exponential domination and
will be reattached using \cref{thm:flat-invisible}.

\subsection{The central integral and the moving window}

Bromwich inversion and the scaling $\theta=v/\sqrt\lambda$ give
\begin{equation}
 \mathcal S(x)=\frac1{\sqrt{2\pi}}
 \int_{\R}\exp\Phi_\lambda(v/\sqrt\lambda)\dd v.
 \label{eq:normalized-saddle-integral}
\end{equation}
The exact phase separates as
\begin{equation}
 \Phi_\lambda(v/\sqrt\lambda)
 =-\frac{v^2}{2}+
 \sum_{m=1}^{M}\varepsilon^mE_m(\lambda,v)
 +R_{M+1}(\lambda,v),
 \qquad
 \varepsilon=\lambda^{-1/2},
 \label{eq:central-exponent-expansion}
\end{equation}
inside an order-dependent central window
\begin{equation}
 |v|\le A_N(\lambda)=\sqrt{32(N+1)\log\lambda}.
 \label{eq:central-window-lambda}
\end{equation}
The all-orders Taylor modules prove, uniformly on this moving window, that the truncated
perturbation tends to zero and that the analytic remainder is dominated by the first omitted
power times a polynomial in $v$.  The slow growth of $A_N$ is crucial:
\begin{equation}
 \lambda^{-1/2}A_N(\lambda)^d\longrightarrow0
 \qquad\text{for every fixed }d.
 \label{eq:epsilon-log-radius-decay}
\end{equation}

\subsection{Finite exponential substitution}

Fix the requested logarithmic order $N$.  Choose the exponent and scalar-exponential
truncation orders high enough that every omitted term is
$O(\varepsilon^{2N})=O(\lambda^{-N})$.  Apply the exact finite-remainder identity
\eqref{eq:finite-remainder-evaluated} pointwise with
\begin{equation}
 x=\varepsilon,
 \qquad
 E_m=E_m(\lambda,v).
 \label{eq:finite-exp-saddle-substitution}
\end{equation}
This yields, on the central window,
\begin{equation}
 \exp\!\left(\sum_{m<M}\varepsilon^mE_m\right)
 =\sum_{k<2N}\varepsilon^ka_k(\lambda,v)
 +O\!\left(\varepsilon^{2N}P_N(v)\right),
 \label{eq:central-reference-expansion}
\end{equation}
where $P_N$ is a fixed-degree polynomial with bounded periodic coefficient families.  The
analytic remainder $R_{M+1}$ is absorbed by a finite scalar exponential estimate because the
whole perturbation is eventually bounded by one on the central window.

\subsection{Gaussian contraction and the disappearance of half-orders}

Multiply \eqref{eq:central-reference-expansion} by $\e^{-v^2/2}$ and integrate.  Parity gives
\begin{equation}
 \Gcal[a_{2j+1}(\lambda,\cdot)]=0.
 \label{eq:odd-saddle-contractions-zero}
\end{equation}
For the even terms,
\begin{equation}
 \frac1{\sqrt{2\pi}}
 \int_\R\e^{-v^2/2}a_{2j}(\lambda,v)\dd v=B_j(\lambda).
 \label{eq:even-saddle-contraction}
\end{equation}
The central-window integral may be replaced by the whole-line integral because
\cref{thm:all-orders-polynomial-tail} makes the complementary Gaussian-polynomial tail
$O(\lambda^{-N})$.

The genuine saddle integrand outside the central window is controlled separately.  The
minor-arc estimates combine strict off-axis decay of the negative-Laplace product with
far-tail polynomial bounds.  They show that the noncentral part of
\eqref{eq:normalized-saddle-integral} is also $O(\lambda^{-N})$ for every fixed $N$.
Consequently
\begin{equation}
 \mathcal S(x)=
 1+\sum_{j=1}^{N-1}\frac{B_j(\lambda)}{\lambda^j}
 +O(\lambda^{-N}).
 \label{eq:full-mass-expansion-supp}
\end{equation}
This is the full mass expansion in the sense of \cref{def:parameter-expansion}.

\subsection{Logarithmic transfer and restoration of the flat tail}

The ratio in \eqref{eq:full-mass-expansion-supp} tends to $1$ and is positive.  Apply the
formal logarithm recurrence and its finite analytic remainder theorem:
\begin{equation}
 \log\mathcal S(x)=
 \sum_{j=1}^{N-1}\frac{A_j(\lambda)}{\lambda^j}
 +O(\lambda^{-N}).
 \label{eq:full-log-mass-expansion-supp}
\end{equation}
Now restore $E(2^\lambda)$ in \eqref{eq:exact-error-separation-supp}.  It is flat by
\eqref{eq:forward-tail-flat-supp}, so \cref{thm:flat-invisible} leaves every coefficient
unchanged.

\begin{theorem}[Full exact-phase logarithmic expansion]
\label{thm:full-exact-phase-expansion}
For every $N\ge1$, as $x\downarrow0$,
\begin{equation}
 \boxed{
 \log F(x)=\mathcal M(x)+
 \sum_{j=1}^{N-1}\frac{A_j(\lambda(x))}{\lambda(x)^j}
 +O\!\left(\lambda(x)^{-N}\right).}
 \label{eq:full-exact-phase-expansion}
\end{equation}
The coefficient functions $A_j$ are bounded, smooth, and one-periodic; $A_1$ and $A_2$ are
given by \eqref{eq:A1-jets} and \eqref{eq:A2-jets}.
\end{theorem}

\begin{proof}
The preceding five subsections establish the expansion on the dyadic logarithmic coordinate.
Use \cref{thm:small-log-expansion-equivalence} to transport it to the small-positive filter.
The exact phase identity $\lambda(2^{-t})$ is retained throughout, so the coefficient
families pull back without approximation.
\end{proof}

Because $\lambda(x)$ is comparable to $-\log x$, the remainder may be weakened to
\begin{equation}
 O\bigl(({-\log x})^{-N}\bigr),
 \label{eq:full-expansion-neglog-rate}
\end{equation}
but the coefficients still remain evaluated at the exact Lambert phase.

\section{Why the exact phase cannot be casually re-expanded}
\label{sec:exact-phase-boundary}

The Lambert phase itself has an all-orders expansion in
$T=-\log x$ and $\log T$.  It is tempting to substitute that expansion into every
$A_j(\lambda(x))$ and claim an ordinary nested-log series.  The current formal theorem does
not do so, for a good reason.

Let $\lambda_N(x)$ be a finite asymptotic approximation to $\lambda(x)$ and write
\begin{equation}
 \Delta_N(x)=\lambda(x)-\lambda_N(x).
 \label{eq:phase-error}
\end{equation}
For a nonconstant periodic coefficient,
\begin{equation}
 A_j(\lambda)-A_j(\lambda_N)
 =A_j'(\lambda_N)\Delta_N+\frac12 A_j''(\xi)\Delta_N^2.
 \label{eq:periodic-phase-Taylor}
\end{equation}
To preserve an $O(\lambda^{-N})$ total remainder, the phase approximation must be accurate
enough after multiplication by every prefactor $\lambda^{-j}$ and after repeated Taylor
composition through all coefficient families.  A statement merely of the form
$\Delta_N=O(\lambda^{-N})$ at one fixed order is not automatically stable under the entire
triangular composition.

\begin{warning}[The formalized boundary]
The theorem proved in \path{FabiusFullAsymptoticExpansion.lean} is
\eqref{eq:full-exact-phase-expansion}: every periodic coefficient is evaluated at the exact
lower-Lambert phase.  A separate branch expands the phase itself to all orders.  What is not
currently claimed is a single all-orders composition theorem that substitutes the phase
series through every oscillatory coefficient while preserving the requested remainder.
This is a missing formal composition layer, not a defect in the exact-phase expansion.
\end{warning}

The exact-phase form is in any case mathematically preferable.  It isolates the implicit
saddle inversion once, avoids secular phase error, and makes phase-locked subsequences
transparent.  For any fixed $u\in\R$, the sequence
\begin{equation}
 x_n=(n+u)2^{-(n+u)}
 \label{eq:phase-locked-sequence-supp}
\end{equation}
has exact phase $\lambda(x_n)=n+u$, so every periodic coefficient is literally constant
along it:
\begin{equation}
 A_j(\lambda(x_n))=A_j(u).
 \label{eq:phase-locked-coefficients}
\end{equation}
This is the cleanest way to separate algebraic decay from periodic oscillation.

\appendix

\section{Executable recurrences in mathematical pseudocode}
\label{app:algorithms}

The following algorithms are not substitutes for the proofs above.  They are compact
implementations of the finite algebra that appears repeatedly in the formal development.
All arrays are zero-indexed.

\begin{algorithm}[Weighted path sum]
\label{alg:path-sum}
The dynamic program below computes $P_w(n)$ without enumerating compositions.  It is the
triangular recurrence proved equivalent to the composition sum in
\cref{thm:triangular-path-solver}.
\end{algorithm}
\begin{lstlisting}[style=pseudo]
function PathSums(N, weight):
    P[0] := 1
    for n := 1 to N:
        P[n] := 0
        for k := 0 to n - 1:
            P[n] := P[n] + P[k] * weight(k, n)
    return P

function FabiusInverseDyadic(N):
    weight(k, n) := 1 / ((2^n - 1) * factorial(n - k + 1))
    a := PathSums(N, weight)
    for n := 0 to N:
        value[n] := 2^(-binomial(n, 2)) * a[n]
    return value
\end{lstlisting}

\begin{algorithm}[Formal exponential coefficients]
\label{alg:expcoeff}
Given $E[1],\ldots,E[N]$ in a commutative $\Q$-algebra, this returns the first $N$
coefficients of $\exp(\sum_{j\ge1}E[j]u^j)$.
\end{algorithm}
\begin{lstlisting}[style=pseudo]
function ExpCoeff(E, N):
    a[0] := 1
    for n := 0 to N - 1:
        s := 0
        for j := 0 to n:
            s := s + (j + 1) * E[j + 1] * a[n - j]
        a[n + 1] := s / (n + 1)
    return a
\end{lstlisting}

\begin{algorithm}[Formal logarithm coefficients]
\label{alg:logcoeff}
The input is a unit series $1+\sum_{j\ge1}a[j]u^j$.  The output satisfies
$\log(1+\sum a[j]u^j)=\sum b[j]u^j$.
\end{algorithm}
\begin{lstlisting}[style=pseudo]
function LogCoeff(a, N):
    b[0] := 0
    for n := 0 to N - 1:
        s := 0
        for j := 0 to n - 1:
            s := s + (n - j) * b[n - j] * a[j + 1]
        b[n + 1] := a[n + 1] - s / (n + 1)
    return b
\end{lstlisting}

\begin{algorithm}[Normalized Gaussian contraction]
\label{alg:gaussian-contraction}
The moment table can be built once to the degree of the input polynomial.
\end{algorithm}
\begin{lstlisting}[style=pseudo]
function GaussianMoments(D):
    mu[0] := 1
    if D >= 1: mu[1] := 0
    for n := 0 to D - 2:
        mu[n + 2] := (n + 1) * mu[n]
    return mu

function GaussianContract(coeff, D):
    mu := GaussianMoments(D)
    result := 0
    for k := 0 to D:
        result := result + coeff[k] * mu[k]
    return result
\end{lstlisting}

\begin{algorithm}[Saddle coefficient pipeline]
\label{alg:saddle-pipeline}
This is the finite algebra at fixed logarithmic order $N$.  The analytic proof is what
justifies applying it to the central integral with a uniform remainder.
\end{algorithm}
\begin{lstlisting}[style=pseudo]
function SaddleCoefficients(N, phase):
    # 1. Build bounded exponent jets Z[0 .. 2N - 1].
    Z := ClosedJets(2*N - 1, phase)

    # 2. Build exponent polynomials E[1 .. 2N].
    for m := 1 to 2*N:
        E[m] := i^m * (Z[m - 1] * X^m / factorial(m)
                       + (-1)^(m + 1) * X^(m + 2) / (m + 2))

    # 3. Exponentiate formally in epsilon.
    a := ExpCoeff(E, 2*N)

    # 4. Contract even epsilon orders against the normalized Gaussian.
    B[0] := 1
    for j := 1 to N:
        B[j] := GaussianContract(coefficients(a[2*j]), degree(a[2*j]))

    # 5. Convert mass coefficients to logarithmic coefficients.
    A := LogCoeff(B, N)
    return (B, A)
\end{lstlisting}

\section{Lean declaration and module concordance}
\label{app:concordance}

File names below are relative to
\begin{center}
 \repo/\texttt{Lean/FabiusFunction/}.
\end{center}
The table lists the principal declarations, not every supporting lemma.

\begingroup\small
\begin{longtable}{@{}>{\raggedright\arraybackslash\scriptsize}p{0.36\textwidth}
                         >{\raggedright\arraybackslash}p{0.57\textwidth}@{}}
\toprule
Module & Content explained in this companion\\
\midrule
\endhead
\path{OriginalCharacterization.lean} &
\lean{IsOriginalFabius.mk\_of\_derivative\_law}: the support/mean-value argument showing
that scale positivity follows from the other original-characterization assumptions;
\cref{sec:scale-positive}.\\
\path{FabiusInverse.lean} &
\lean{fabiusOrderIso}, \lean{fabiusInv}, \lean{fabiusReal\_fabiusInv},
\lean{fabiusInv\_fabiusReal}, \lean{continuous\_fabiusInv},
\lean{strictMonoOn\_fabiusInv}, \lean{fabiusInv\_one\_sub},
\lean{fabiusInv\_fabiusDyadicUnit}, \lean{fabiusInv\_hasDerivAt},
\lean{deriv\_fabiusInv}, \lean{deriv\_fabiusInv\_eq\_inv\_two\_mul\_rvachevUp},
\lean{deriv\_fabiusInv\_pos}, \lean{fabiusInv\_contDiffOn\_Ioo},
\lean{deriv\_deriv\_fabiusInv}, \lean{strictConcaveOn\_fabiusInv\_firstHalf},
\lean{strictConvexOn\_fabiusInv\_secondHalf}, the transported root-free bounds,
\lean{mul\_lt\_fabiusInv}, \lean{tendsto\_fabiusInv\_div\_atTop}, and reflected companions;
\cref{sec:inverse,sec:inverse-steepness}.\\
\path{ElementaryFunction.lean} &
The inductive predicate \lean{IsElementary}, composition closure, derived elementary
operations, analytic loci, the finite-bad-value composition theorem, and the dense-open
analytic-locus theorem; \cref{sec:elementary}.\\
\path{NotElementary.lean} &
No elementary function agrees with a bounded Fabius function, $\Up$, or the signed extension
on a set with nonempty interior; \cref{sec:not-elementary}.\\
\path{ProbabilityRepresentation.lean} &
The product sample space, uniformly convergent coordinate sum, pushforward law, head--tail
split, coordinate reflection, and identification of the CDF with $F$;
\cref{sec:product-law}.\\
\path{ProbabilityLaplaceMoments.lean} &
Almost-sure support, restriction to $[0,1]$, survival function, the compact-support Fubini
identity, raw and endpoint moment identifications, zero-tilt normalization, and positivity;
\cref{sec:survival-fubini,sec:probability-laplace}.\\
\path{LaplaceMoments.lean} &
\lean{tiltedSurvivalMoment}, \lean{fabiusLaplaceMoment}, the differential recurrence,
iterated derivatives, smoothness, and evaluation at zero tilt; \cref{sec:probability-laplace}.\\
\path{StepMeasureBridge.lean} &
Almost-everywhere replacement of half-endpoint indicators, exact overlap integrals, atomic
measure comparisons, the shrunken/enlarged interval sandwich, portmanteau, and convergence of
every interval integral; \cref{sec:step-measure-bridge}.\\
\path{FabiusInverseDyadicClosedForm.lean} &
Compositions, path weights, last-edge decomposition, generic triangular solver, uniqueness,
and the exact composition formula for $F(2^{-n})$; \cref{sec:path-sums,sec:fabius-composition-formula}.\\
\path{FabiusDiscreteLimitToeplitz.lean} &
Corrected prefix length, exact $q$-binomial row mass, total-variation bound $16$, the generic
finite-row Toeplitz theorem, and the specialization $q(n,j)=2n-j$;
\cref{sec:toeplitz,sec:fabius-toeplitz}.\\
\path{SaddleExpansionAlgebra.lean} &
\lean{expCoeff}, its differential equation, uniqueness, naturality, rescaling, addition,
parity, and \lean{HasAsymptoticExpansion}; \cref{sec:poincare-api,sec:exp-coeff}.\\
\path{SaddleLogExpansionAlgebra.lean} &
\lean{logCoeff}, the identity $AB'=A'$, low-order formulas, uniqueness, naturality,
rescaling, and parity; \cref{sec:log-coeff}.\\
\path{SaddleExpansionFiniteRemainder.lean} &
Finite exponent substitution, equality of coefficients below the truncation order,
divisibility of the defect by $X^L$, and the evaluated exact remainder;
\cref{sec:finite-remainder}.\\
\path{SaddleExpansionSmallArgument.lean} &
Exact transfer of full expansions between $t\to\infty$ and $x=2^{-t}\downarrow0$;
\cref{sec:poincare-api}.\\
\path{SaddleExpansionFlat.lean} &
Zero-coefficient expansion of all-orders-flat errors and invariance of every coefficient under
addition or subtraction of such an error; \cref{sec:poincare-api,sec:full-assembly}.\\
\path{GaussianPolynomialContraction.lean} &
Gaussian integrability, odd moments, integration-by-parts recurrence, normalized moments,
linear polynomial contraction, and the whole-line integral identity;
\cref{sec:gaussian-contraction}.\\
\path{GaussianPolynomialTail.lean} &
The factorial coefficient weight and explicit complementary-interval tail bound;
\cref{sec:gaussian-tails}.\\
\path{GaussianPolynomialTailAllOrders.lean} &
The radius $\sqrt{32(N+1)\log b}$ and arbitrary algebraic tail order for bounded polynomial
families; \cref{sec:gaussian-tails}.\\
\path{FabiusSaddleJetClosedForm.lean} &
Closed formula for the periodic and bounded exponent jets, including harmonic constants and
the shifted falling-factorial coefficients; \cref{sec:saddle-jets}.\\
\path{FabiusSaddleJetStirling.lean} &
Identification $c(n,m)=s(n+1,m+1)$ with signed Stirling numbers of the first kind;
\cref{sec:saddle-jets}.\\
\path{FabiusSaddleExponentClosedForm.lean} &
Cancellation of factorials and the extra power of $\ii$, yielding the two-monomial formula
\eqref{eq:closed-exponent-coefficient}; \cref{sec:saddle-jets}.\\
\path{FabiusSaddleLeadingCoefficient.lean} &
The unique newest-jet monomial and coefficient $(-1)^j/(2^jj!)$ in the $j$th mass
coefficient; \cref{sec:saddle-coeff-engine}.\\
\path{FabiusSecondSaddleCorrection.lean} &
Low-order formal exponential identities, explicit order-four polynomial, Gaussian contraction,
closed $B_2$, and closed $A_2$; \cref{sec:second-correction}.\\
\path{FabiusSecondSaddleExpansion.lean} &
The explicit expansion through $A_2/\lambda^2$ with $O(\lambda^{-3})$ and the weakened
$O(({-\log x})^{-3})$ form; \cref{sec:second-correction,sec:full-assembly}.\\
\path{FabiusSaddleMassAllOrders.lean} &
Finite exponential substitution, bounded jet families, central reference polynomial,
Gaussian whole-line contraction, and all-orders normalized saddle mass;
\cref{sec:full-assembly}.\\
\path{FabiusFullAsymptoticExpansion.lean} &
Transfer from complex kernel mass to the real saddle ratio, logarithmic coefficient transfer,
addition of the flat forward tail, exact-phase full expansion, and small-argument transport;
\cref{sec:full-assembly,sec:exact-phase-boundary}.\\
\bottomrule
\end{longtable}
\endgroup

\section{Audit matrix against the main article}
\label{app:audit-matrix}

\begingroup\small
\begin{longtable}{@{}p{0.28\textwidth}p{0.18\textwidth}p{0.46\textwidth}@{}}
\toprule
Topic & Status before this companion & Added here\\
\midrule
\endhead
Original scale positivity & Assumed in the source statement & Mean-value proof that the dilation
parameter is automatically positive before its exact value $k=2$ is derived.\\
Total inverse $I$ & Absent & Global clamped construction, order calculus, reflection, exact
dyadic inverse values, interior calculus consequences, endpoint steepness, and non-H\"older
corollaries.\\
Elementary function class & Absent & Exact inductive closure, totalization convention,
finite-bad-value composition, and dense-open analytic locus.\\
Non-elementarity & Absent & No local agreement on any set with nonempty interior for $F$,
$\Up$, and the signed extension.\\
Product probability law & Headline proof only & Uniform convergence and measurability,
head--tail pushforward identity, coordinate reflection, and global support bookkeeping.\\
Probability--Laplace bridge & Absent & Survival-function Fubini identity, identification of
raw, endpoint, and tilted moments, complete monotonicity, and half-moment interpretation.\\
Atomic-to-step bridge & Compressed & Half-endpoint null-set conversion, exact cell overlap,
shrunken/enlarged interval sandwich, portmanteau, and moving-boundary squeeze.\\
Weighted path solver & Explicitly omitted & Generic semiring theorem, uniqueness, edge-count
decomposition, Fabius specialization, and worked values.\\
Finite-row Toeplitz theorem & Explicitly omitted & Abstract normed-ring proof, exact Fabius
row mass, uniform variation bound, tail-index condition, and corrected prefix rounding.\\
Parameter-dependent full expansions & Only summarized & Exact API, coefficient boundedness,
closure, flat errors, and coordinate transfer.\\
Formal exponential/logarithm & Explicitly omitted & Recurrences, formal differential equations,
uniqueness, low orders, naturality, rescaling, addition, and parity.\\
Finite formal-to-analytic remainder & Absent & Exact divisibility by the first omitted power
and evaluated polynomial remainder.\\
Gaussian contraction/tails & Explicitly omitted & Moment recurrence, contraction functional,
factorial coefficient weight, explicit tail, and all-orders radius.\\
Closed saddle jets & Stated, lightly motivated & Operator-product solution, harmonic constants,
shifted signed Stirling identification, and first four bounded jets.\\
Second saddle correction & Not yet in pinned snapshot & Complete $E_1$--$E_4$ expansion,
order-four polynomial, Gaussian contraction, $B_2$, and closed $A_2$.\\
Full saddle assembly & Proof architecture only & Exact sequence of finite remainder, central
window, contraction, Gaussian tail, minor arc, log transfer, flat-tail restoration, and filter
transport.\\
Exact phase boundary & Warning only & Quantified explanation of why all-orders substitution
through oscillatory coefficients needs a separate composition theorem.\\
\bottomrule
\end{longtable}
\endgroup

The following substantial topics are intentionally not repeated because the main article
already treats them at proof level: construction and uniqueness of $F$ and $\Up$; the original
Rvachev characterization apart from the redundancy isolated in \cref{sec:scale-positive}; the infinite sinc product and Fourier decay; support, positivity,
shape, and nowhere analyticity; exact moment recurrences and denominator arithmetic; dyadic
rational evaluation, $q$-Pochhammer and $q$-binomial formulas; Thue--Morse interpolation and
iterated difference identities; Legendre expansions and least-squares approximation; and the
computational API.  Their use as input above is always signposted.

\section{Notation crosswalk}
\label{app:notation}

\begin{longtable}{@{}p{0.18\textwidth}p{0.29\textwidth}>{\raggedright\arraybackslash\footnotesize}p{0.43\textwidth}@{}}
\toprule
This companion & Main-article notation & Principal Lean declaration\\
\midrule
\endhead
$I(y)$ & inverse of bounded $F$ & \lean{fabiusInv}\\
$\mu$ & law of the weighted uniform sum & \lean{weightedSumDistribution}\\
$J_k(s)$ & tilted survival moment & \lean{tiltedSurvivalMoment}\\
$M_k(s)$ & tilted probability/Fabius moment & \lean{fabiusLaplaceMoment}\\
$d_n$ & exact half moment & \lean{halfMoment}\\
$P_w(n)$ & weighted composition/path sum & generic path-sum definition in
\path{FabiusInverseDyadicClosedForm.lean}\\
$\sigma$ & abstract small scale & first argument of \lean{HasAsymptoticExpansion}\\
$a_n$ & formal exponential coefficient & \lean{SaddleExpansion.expCoeff}\\
$b_n$ & formal logarithm coefficient & \lean{SaddleExpansion.logCoeff}\\
$\Gcal[p]$ & normalized Gaussian contraction &
\lean{SaddleExpansion.}\newline\lean{gaussianPolynomialContraction}\\
$Z_n(t)$ & bounded exponent jet $\widetilde{\mathcal Z}_n$ &
\lean{negativeLaplaceBoundedExponentJet}\\
$B_j(t)$ & normalized saddle mass coefficient & \lean{fabiusSaddleMassCoefficient}\\
$A_j(t)$ & logarithmic saddle coefficient & \lean{fabiusSaddleLogCoefficient}\\
$\lambda(x)$ & exact lower-Lambert phase & \lean{fabiusLambertPhase}\\
$\mathcal M(x)$ & corrected sharp Lambert main & \lean{fabiusSharpLambertMain}\\
$\mathcal S(x)$ & normalized real saddle mass/ratio & \lean{fabiusSaddleRatio} and the
kernel-mass bridge\\
\bottomrule
\end{longtable}

\section{A compact dependency diagram}
\label{app:dependency}

The proof dependencies of the material unique to this companion can be read from the
following diagram.  An arrow means that the target uses the source as a mathematical input,
not merely as a Lean import.

\begin{center}
\begin{minipage}{0.96\textwidth}\footnotesize
\begin{align*}
 &\text{support + interior positivity + mean-value theorem + derivative law}
   \longrightarrow k>0,\\[0.35ex]
 &\text{strict monotonicity + surjectivity + flatness}
   \longrightarrow \text{total inverse}
   \longrightarrow \text{endpoint steepness},\\[0.35ex]
 &\text{nowhere analyticity + dense elementary analytic locus}
   \longrightarrow \text{strong non-elementarity},\\[0.35ex]
 &\text{product law}
   \longrightarrow \text{survival identity}
   \longrightarrow \text{Laplace/raw/endpoint moment bridge},\\[0.35ex]
 &\text{weak convergence + null boundaries + cell geometry}
   \longrightarrow \text{step-integral convergence},\\[0.35ex]
 &\text{triangular recurrence}
   \longrightarrow \text{weighted paths}
   \longrightarrow \text{composition formula},\\[0.35ex]
 &\text{$q$-binomial row mass + variation + tail indices}
   \longrightarrow \text{discrete Toeplitz limit},\\[0.35ex]
 &\text{closed jets}
   \longrightarrow \text{exponent coefficients}
   \longrightarrow \text{formal exponential}\\[-0.1ex]
 &\hspace{8em}\longrightarrow \text{Gaussian contraction}
   \longrightarrow \text{mass coefficients},\\[0.35ex]
 &\text{mass coefficients}
   \longrightarrow \text{formal logarithm}
   \longrightarrow \text{log coefficients},\\[0.35ex]
 &\text{finite remainders + central Taylor + Gaussian tails + minor arcs}\\[-0.1ex]
 &\hspace{8em}\longrightarrow \text{all-orders saddle mass},\\[0.35ex]
 &\text{all-orders mass + exact reduction + flat tail + coordinate transfer}\\[-0.1ex]
 &\hspace{8em}\longrightarrow \text{full exact-phase expansion}.
\end{align*}
\end{minipage}
\end{center}

\nocite{fabius1966,arias1982,arias05442,arias2018,mainarticle,repo}
\begin{thebibliography}{99}

\bibitem{fabius1966}
J.~Fabius,
\newblock A probabilistic example of a nowhere analytic $C^\infty$-function,
\newblock \emph{Zeitschrift f\"ur Wahrscheinlichkeitstheorie und Verwandte Gebiete}
\textbf{5} (1966), 173--174.

\bibitem{arias1982}
J.~Arias de Reyna,
\newblock On the Fabius function,
\newblock \emph{The American Mathematical Monthly} \textbf{89} (1982), 291--298.

\bibitem{arias05442}
J.~Arias de Reyna,
\newblock An infinitely differentiable function with compact support: Definition and properties,
\newblock arXiv:1702.05442.

\bibitem{arias2018}
J.~Arias de Reyna,
\newblock The asymptotic behavior of the Fabius function,
\newblock arXiv:1702.06487, version 3, 2018.

\bibitem{mainarticle}
V.~Reshetnikov,
\newblock \emph{The Fabius Function and Rvachev's Up-Function: Exact Arithmetic,
Probability, Fourier Analysis, and Corrected Sharp Asymptotics},
\newblock \texttt{ProveIt},
\path{Analysis/FabiusFunction/docs/Fabius_Function_and_Rvachev_Up/}.

\bibitem{repo}
V.~Reshetnikov,
\newblock \texttt{ProveIt}: formal development of the Fabius function and Rvachev's
up-function,
\newblock GitHub repository, snapshot
\texttt{5e80805b1ee92697723949d125aa0d6dbf32f538}, 25 August 2026,
\newblock \url{https://github.com/VladimirReshetnikov/ProveIt/tree/main/Analysis/FabiusFunction}.

\end{thebibliography}

\end{document}
